\chapter{Компиляторы схем}\label{chap:circuit-compilers}
% finite field arithmetics https://johnkerl.org/doc/ffcomp.pdf
Как было показано ранее в главе~\ref{sec:statements}, утверждения в формальных языках могут быть формализованы как утверждения о принадлежности или знании, при этом алгебраические схемы (algebraic circuits) и системы ограничений ранга~1 (Rank-1 Constraint Systems) являются двумя ключевыми средствами определения таких языков.

Однако как алгебраические схемы, так и системы ограничений ранга~1 существенно отличаются от традиционных парадигм программирования и создают трудности для разработчиков. Задача написания реальных приложений в виде схем и проверки их через системы ограничений ранга~1 столь же сложна, как написание кода на низкоуровневых языках, таких как ассемблер (assembly). Для упрощения проектирования сложных утверждений необходим фреймворк компилятора (compiler), который может переводить языки высокого уровня (high-level languages) в арифметические схемы и связанные с ними системы ограничений ранга~1.

Эти языки программирования позволяют более интуитивно и эффективно создавать и тестировать арифметические схемы, освобождая разработчиков от сложностей представлений R1CS. В результате они могут сосредоточиться на логике схем, которые они собираются построить, в то время как компилятор (compiler) берёт на себя всё остальное.

Дополнительным преимуществом языков программирования, компилируемых в R1CS, является то, что они могут генерировать не только представления R1CS, но и другие программы, способные эффективно вычислять значения присваиваний (assignment) схемы. Это облегчает интеграцию схем в различные платформы и среды без необходимости учитывать детали нижележащей реализации.

Как было показано в главе~\ref{sec:statements} в разделах~\ref{sec:R1CS} и~\ref{sec:circuits}, как арифметические схемы, так и системы ограничений ранга~1 обладают свойством модульности (modularity) \ref{sec:R1CS_modularity}, что позволяет синтезировать сложные схемы из более простых. Многие компиляторы схем/R1CS придерживаются базового подхода, предоставляя библиотеку атомарных и простых схем вместе со средствами объединения этих строительных блоков в более сложные системы.

Эта глава даёт обзор базовых концепций в области компиляторов схем (circuit compilers), представляя игрушечный язык, который можно вручную ``скомпилировать'' в графические представления алгебраических схем (algebraic circuits) и связанных с ними систем ограничений ранга~1. Затем в главе рассматриваются реальные компиляторы и языки более высокого уровня, которые они поддерживают.

Глава начинается с общего введения в игрушечный язык программирования и в реальные языки, за которым следует обсуждение атомарных типов, таких как булевы значения (booleans) и целые числа без знака (unsigned integers). Затем определяются примитивы управления потоком выполнения, включая условный оператор if-then-else и ограниченный цикл. Глава завершается обзором базовых функциональных примитивов, таких как криптография на эллиптических кривых, обычно называемых в литературе ``гаджетами'' (gadgets).

\section{Язык ``ручки и бумаги''} Чтобы объяснить базовые концепции компиляторов схем (circuit compilers) и связанных с ними языков высокого уровня, мы выведем неформальный игрушечный язык и связанный с ним ``мозговой'' компилятор (brain-compiler), который назовём \lgname{PAPER} (\textbf{P}en-\textbf{A}nd-\textbf{P}aper \textbf{E}xecution \textbf{R}ules). \lgname{PAPER} позволяет программистам определять утверждения в псевдокоде, похожем на Rust. Язык вдохновлён \lgname{zokrates} и \lgname{circom}.
 
\subsection{Грамматика}
В \lgname{PAPER} любое утверждение определяется как упорядоченный список функций, где любая функция должна быть объявлена в списке до того, как она будет вызвана в другой функции этого списка. Последняя запись в утверждении должна быть специальной функцией, называемой \texttt{main}. Функции принимают список типизированных параметров в качестве входных данных и вычисляют кортеж типизированных переменных (variables) в качестве выходных, где type\_functions --- это специальные функции, определяющие, как преобразовать один тип в другой, в конечном итоге преобразуя любой тип в элементы базового поля, над которым определена схема.

Любое утверждение параметризовано полем, над которым будет определена схема, и имеет дополнительные необязательные параметры беззнакового типа, необходимые для определения размера массивов или счётчика ограниченных циклов. Следующее определение делает грамматику утверждения точной с помощью описания в стиле командной строки:
% used the command line style formalized e.g. here: http://docopt.org/
\begin{lstlisting}
statement <Name> {F:<Field> [ , <N_1: unsigned>,... ] } {
  [fn <Name>([[pub]<Arg>:<Type>,...]) -> (<Type>,...){
    [let [pub] <Var>:<Type> ;... ]
    [let const <Const>:<Type>=<Value> ;... ]
    Var<==(fn([<Arg>|<Const>|<Var>,...])|(<Arg>|<Const>|<Var>)) ;
    return (<Var>,...) ;
  } ;...]
  fn main([[pub]<Arg>:<Type>,...]) -> (<Type>,...){
    [let [pub] <Var>:<Type> ;... ]
    [let const <Const>:<Type>=<Value> ;... ]
    Var<==(fn([<Arg>|<Const>|<Var>,...])|(<Arg>|<Const>|<Var>)) ;
    return (<Var>,...) ;
  } ;
}
\end{lstlisting}
Аргументы функций и переменные (variables) по умолчанию являются witness-переменными, но могут быть объявлены как instance с помощью спецификатора \texttt{pub}. Объявление аргументов и переменных как instance всегда перезаписывает любые предыдущие или конфликтующие witness-объявления. Каждый аргумент, константа (constant) или переменная (variable) имеет тип, и каждый тип определяется как функция, преобразующая этот тип в другой тип. Для того чтобы программа на \lgname{PAPER} успешно скомпилировалась, все преобразования типов должны быть составлены таким образом, чтобы конечный тип был базовым полем, над которым определена схема:
\begin{lstlisting}
type_function <TYPE>( t1 : <TYPE_1>) -> TYPE_2{
  let t2: TYPE_2 <== fn(TYPE_1)
  return t2
}
\end{lstlisting}
Многие реальные языки схем основаны на подобном, но, конечно, более изощренном подходе, чем \lgname{PAPER}. Цель \lgname{PAPER} --- показать базовые принципы компиляторов схем (circuit compilers) и связанных с ними языков высокого уровня (high-level languages).
\begin{example}Чтобы лучше понять грамматику \lgname{PAPER}, следующий код является корректным высокоуровневым кодом, соответствующим грамматике языка \lgname{PAPER}, при условии, что все типы в этом коде определены где-либо ещё.
\begin{lstlisting}
statement MOCK_CODE {F: F_43, N_1 = 1024, N_2 = 8} {
  fn foo(in_1 : F, pub in_2 : TYPE_2) -> F {
    let const c_1 : F = 0 ;
    let const c_2 : TYPE_2 = SOME_VALUE ;
    let pub out_1 : F ;
    out_1<== c_1 ;
    return out_1 ;
  } ;
  
  fn bar(pub in_1 : F) -> F {
    let out_1 : F ;
    out_1<==foo(in_1);
    return out_1 ;
  } ;
    
  fn main(in_1 : TYPE_1)->(F, TYPE_2){
    let const c_1 : TYPE_1  = SOME_VALUE ;
    let const c_2 : F = 2;
    let const c_3 : TYPE_2  = SOME_VALUE  ;
    let pub out_1 : F ;
    let out_2 : TYPE_2 ;
    c_1 <== in_1 ;
    out_1 <== foo(c_2) ;
    out_2 <== TYPE_2 ;
    return (out_1,out_2) ;
  } ;
}
\end{lstlisting}
\end{example}
\subsection{Фазы выполнения} В отличие от обычных исполняемых программ, программы для компиляторов схем (circuit compilers) имеют два режима выполнения. Первый режим, обычно называемый \term{setup phase} (фаза настройки), выполняется для генерации схемы (circuit) и связанной с ней системы ограничений ранга~1 (Rank-1 Constraint System), которая затем обычно используется как вход для некоторой системы доказательств с нулевым разглашением (zero-knowledge proof system), как объясняется в~\ref{chapter:zk-protocols}.

Второй режим выполнения обычно называется \term{prover phase} (фаза доказывающего). На этом этапе некоторое присваивание (assignment) всем instance-переменным схемы обычно задаётся как вход, и задача доказывающего (prover) состоит в том, чтобы вычислить корректное присваивание (assignment) всем witness-переменным схемы. В зависимости от случая использования, это корректное присваивание затем либо напрямую используется как конструктивное доказательство правильного выполнения схемы, либо передаётся как вход в алгоритм генерации доказательства некоторой системы доказательств с нулевым разглашением, где полноразмерное, не скрывающее конструктивное доказательство обрабатывается в краткое доказательство с различными уровнями нулевого разглашения.

Современные языки схем и связанные с ними компиляторы (compilers) абстрагируют эти две фазы и предоставляют разработчику единый \term{interphase} (интерфейс), который затем пишет единую программу, которую можно использовать в обеих фазах.

Чтобы дать читателю чёткое концептуальное различие между двумя фазами, \lgname{PAPER} сохраняет их раздельными. Код на \lgname{PAPER} может быть ``мозговым компилятором'' (brain-compiled) во время \term{setup-phase} (фазы настройки) в подходе ``ручка и бумага'' в графическое представление схемы. Как только схема выведена, она может быть выполнена в \term{prover phase} (фазе доказывающего) для генерации корректного присваивания (assignment). Корректное присваивание затем интерпретируется как конструктивное доказательство знания в связанном языке.
\paragraph{Фаза настройки} В \lgname{PAPER} задача фазы настройки (setup phase) состоит в компиляции кода на языке \lgname{PAPER} в визуальное представление алгебраической схемы (algebraic circuit). Выведение схемы из кода в стиле ``ручка и бумага'' --- это то, что мы называем \term{brain-compiling} (``мозговая'' компиляция).

Дано некоторое описание утверждения, соответствующее правильной грамматике, мы начинаем процесс графической компиляции схемы с пустой схемы, сначала компилируем функцию main, а затем индуктивно компилируем все остальные функции по мере их вызова в процессе.

Для каждой функции, которую мы в данный момент компилируем, мы рисуем прямоугольный узел для каждого входного аргумента, каждой переменной (variable) и каждой константы (constant) этой функции. Если узел представляет переменную (variable), мы помечаем его именем этой переменной, а если он представляет константу (constant), мы помечаем его значением этой константы. Мы группируем аргументы в подграф, помеченный ``inputs'' (входы), и возвращаемые значения в подграф, помеченный ``outputs'' (выходы). Затем мы группируем всё в подграф и помечаем этот подграф именем функции.

После этого мы должны выполнить проверку согласованности и типов для каждого вхождения оператора присваивания (assignment) \texttt{<==}. Мы должны убедиться, что выражение в правой части оператора корректно определено и что типы обеих сторон совпадают.

Затем мы компилируем правую сторону каждого вхождения оператора присваивания (assignment) \texttt{<==}. Если правая сторона является константой или переменной (variable), определённой в этой функции, мы рисуем пунктирную линию от прямоугольного узла, представляющего правую сторону \texttt{<==}, к прямоугольному узлу, представляющему левую сторону того же оператора. Если правая сторона представляет аргумент этой функции, мы рисуем линию от прямоугольного узла, представляющего левую сторону \texttt{<==}, к прямоугольному узлу, представляющему правую сторону того же оператора.

Если правая сторона оператора \texttt{<==} является функцией, мы заглядываем в нашу базу данных, находим связанную с ней схему и рисуем её. Если с этой функцией ещё не связана ни одна схема, мы повторяем процесс компиляции для этой функции, рисуя рёбра от аргументов функции к её входным узлам и от выходных узлов функции к узлам в правой части \texttt{<==}.

В процессе этого метки рёбер рисуются в соответствии с определяющими правилами алгебраических схем из~\ref{def:algebraic-circuit}. Если связанная переменная представляет witness-переменную, мы используем метку $W$ для обозначения witness, а если она представляет instance-переменную, мы используем метку $I$ для обозначения instance. Переменные (variables) по умолчанию являются witness, а спецификатор $pub$ указывает, что переменная является instance.

После завершения этого мы компилируем все встречающиеся типы всех переменных в функции, компилируя type\_function каждого типа. Мы делаем это индуктивно, пока не достигнем типа базового поля. Схемы не имеют понятия типов, только элементов поля; следовательно, каждый тип должен быть скомпилирован в полевой тип в последовательности шагов компиляции.

Компиляция останавливается, как только мы индуктивно заменили все функции их схемами. Результат --- это схема, содержащая множество ненужных прямоугольных узлов. На заключительном шаге \term{optimization} (оптимизации) все прямоугольные узлы, напрямую связанные друг с другом, сворачиваются в один узел, а все прямоугольные узлы, представляющие одинаковые константы (constants), сворачиваются в один узел.

Конечно, ``мозговой'' компилятор (brain-compiler) \lgname{PAPER} не определён должным образом в формальном смысле. Его цель --- выделить важные шаги, которые проходят реальные компиляторы (compilers) в своих фазах настройки (setup phases).
\begin{example}[Тривиальная схема]\label{ex:trivial-circuit} Чтобы дать интуицию о том, как писать и компилировать схемы на языке \lgname{PAPER}, рассмотрим следующее описание утверждения:
\begin{lstlisting}
statement trivial_circuit {F:F_13} {
  fn main{F}(in1 : F, pub in2 : F) -> (F,F){
    let const outc1 : F = 0 ;   
    let const inc1 : F = 7 ;
    let out1 : F ;
    let out2 : F ;
    out1 <== inc1;
    out2 <== in1;
    outc1 <== in2;
    return (out1, out2) ;
  }
}
\end{lstlisting} 
Чтобы ``мозговым компилятором'' (brain-compile) скомпилировать это утверждение в алгебраическую схему (algebraic circuit) с помощью \lgname{PAPER}, мы начинаем с пустой схемы и вычисляем функцию \texttt{main}, которая является единственной функцией в этом утверждении.

Мы рисуем прямоугольные узлы для каждого аргумента, каждой константы (constant) и каждой переменной (variable) функции и помечаем их их именами или значениями соответственно. Затем мы выполняем проверку согласованности и типов для каждого оператора \texttt{<==} в функции. Поскольку схема просто соединяет входы с выходами и все элементы имеют одинаковый тип, проверка корректна.

Затем мы вычисляем правую сторону операторов присваивания (assignment). Поскольку в нашем случае правая сторона каждого оператора не является функцией, мы рисуем рёбра от прямоугольных узлов в правой части к связанному прямоугольному узлу в левой части. Чтобы пометить эти рёбра, мы используем общие правила алгебраических схем, как определено в~\ref{def:algebraic-circuit}. Согласно этим правилам, каждое входящее ребро стокового узла имеет метку и каждое исходящее ребро от источникового узла имеет метку, если узел помечен переменной. Поскольку узлы, представляющие константы (constants), неявно предполагаются частными, а спецификатор public определяет, помечено ли ребро $W$ или $I$, мы получаем следующую схему:
\begin{center}
\digraph[scale=0.4]{TRIVIAL1}{
	forcelabels=true;
	center=true;
	splines=ortho;
	nodesep= 2.0;
	subgraph clusterINPUTS {
		n1 [shape=box, label="in_1"];
	  n2 [shape=box, label="in_2 = 0"];
	  color= lightgray ;
	  label = "inputs" ;
	}
	
	subgraph clusterOUTPUTS {
		n4 [shape=box, label="out_1"];
	  n5 [shape=box, label="out_2"];
		color= lightgray ;
	  label = "outputs" ;
	}

  subgraph clusterINNER {
  	n3 [shape=box, label="7"];
	  n6 [shape=box, label="0"];
	  color= white ;
  }
	n1 -> n5 [xlabel="W_1"] ;
	n2 -> n6 [xlabel="0"] ;
	n3 -> n4 ;
}
\end{center}
\end{example}
\paragraph{Фаза доказывающего} В \lgname{PAPER} так называемая \term{prover phase} (фаза доказывающего) может быть выполнена, как только фаза настройки (setup phase) сгенерировала графическое представление схемы из связанного с ней высокоуровневого кода. Это делается путём выполнения схемы с присваиванием правильных значений всем входным узлам схемы. Однако, в отличие от большинства реальных компиляторов, \lgname{PAPER} не сообщает доказывающему (prover), как найти правильные входные значения для данной схемы. Реальные языки программирования обычно предоставляют эти данные посредством вычислений, выполняемых вне схемы.
\begin{example} Рассмотрим схему из~\examplename{}~\ref{ex:TJJ-circuit_2}. Корректные присваивания (assignments) к этой схеме являются конструктивными доказательствами того, что пара входов $<I_1,I_2>$ является точкой на кривой tiny-jubjub. Однако схема не предоставляет способа фактически вычислить правильные значения для $I_1$ и $I_2$. Поэтому любая реальная система нуждается во вспомогательном вычислении, предоставляющем эти значения.
\end{example}
\section{Реальные языки схем}
Языки программирования, компилируемые в системы ограничений ранга~1, становятся всё более популярными в областях криптографии и блокчейна. Эти языки предоставляют абстракцию более высокого уровня для проектирования и реализации арифметических схем, которые являются фундаментальными строительными блоками в системах доказательств с нулевым разглашением. Используя эти языки программирования, разработчики могут сосредоточиться на логике схем, которые они хотят построить, в то время как компилятор (compiler) берёт на себя низкоуровневые детали представлений R1CS. Кроме того, эти языки могут выводить не только представление R1CS, но и другие программы, способные эффективно вычислять все значения корректного присваивания (assignment) схемы.

\subsection{Circom}
Circom --- это предметно-ориентированный язык программирования для проектирования арифметических схем. Он используется для построения схем, которые могут быть скомпилированы в системы ограничений ранга~1 и выведены как программы WebAssembly и C++ для эффективного вычисления.

В этом разделе мы приведём примеры написания базовых схем в Circom. Мы будем использовать эти примеры позже для вычисления связанных доказательств в snarkjs.

Чтобы понять Circom, мы сначала должны дать определения терминам \hilight{signals} (сигналы), \hilight{templates} (шаблоны) и \hilight{components} (компоненты), чтобы облегчить лучшее понимание обсуждаемых примеров.

\term{Signal} (сигнал) относится к элементу в лежащем в основе конечном поле $\F$ схемы. Арифметические схемы, созданные с помощью Circom, оперируют сигналами (signals), которые неизменяемы и могут быть определены как входы или выходы. Входные сигналы являются приватными, если не указано иное как public, и все выходные сигналы являются публично доступными. Остальные сигналы являются приватными и не могут быть сделаны публичными. Публичные сигналы являются частью instance, а приватные сигналы --- частью witness в любом корректном присваивании (assignment) схемы.

\term{Template} (шаблон) --- это алгоритм, создающий обобщённые схемы в Circom. Шаблон --- это новый объект схемы, который может использоваться для построения других схем.

\term{Components} (компоненты) определяют арифметическую схему, получая входные сигналы и производя выходные сигналы и промежуточные сигналы, а также набор ограничений (constraints). Компоненты, как и сигналы, также неизменяемы.

\begin{example}
\label{ex:trivial-circuit-circom}
В качестве демонстрации генерации схем на языке Circom мы вновь рассмотрим тривиальную схему из примера~\ref{ex:trivial-circuit}, ранее определённую в \lgname{PAPER}. Следующий код предоставляет одно представление этой схемы на языке Circom:

\begin{lstlisting}
template trivial_circuit() {

    signal input in1 ;
    signal input in2 ;

    var outc1 = 0 ;
    var inc1 = 7 ;
    
    signal output out1 ;
    signal output out2 ;

    out2 <== in1 ;
    outc1 === in2 ;
}

component main = trivial_circuit() ;
\end{lstlisting}

Следует отметить, что лежащее в основе поле не указывается явно в Circom, поэтому константы $0$ и $7$ не имеют непосредственного значения. Для компиляции языка Circom в представление R1CS код должен быть сохранён в текстовый файл. Мы будем называть его \fname{trivial\_circuit.circom}. Этот файл затем может быть скомпилирован с помощью следующей команды:
\begin{lstlisting}
circom trivial_circuit.circom --r1cs --wasm --sym
\end{lstlisting}

Компилятор Circom генерирует три различных файла. Первый файл, \fname{trivial\_circuit.r1cs}, содержит систему ограничений R1CS схемы в пользовательском бинарном формате. Второй файл, \fname{trivial\_circuit.wasm}, является wasm-кодом, который вычисляет witness из данного instance, тем самым давая решение для R1CS. Наконец, третий файл, \fname{trivial\_circuit.sym}, является текстовым файлом символов, необходимым для отладки или аннотированного вывода системы ограничений.
\end{example}

\begin{example}[Задача 3-факторизации в Circom]
\label{ex:3-fac-circom}
В этом примере мы реализуем задачу 3-факторизации \ref{ex:3-factorization} на языке Circom и компилируем в R1CS и генератор утверждения. Чтобы показать, как Circom обрабатывает модульность (modularity) \ref{modularity}, мы напишем код следующим образом:
\begin{lstlisting}
template Multiplier() {
    signal input a ;
    signal input b ;
    signal output c ;
    c <== a*b ;
}


template three_fac () {
    signal input x1 ;
    signal input x2 ;
    signal input x3 ;
    signal output x4 ;
    component mult1 = Multiplier() ;
    component mult2 = Multiplier() ;
    mult1.a <== x1 ;
    mult1.b <== x2 ;
    mult2.a <== mult1.c ;
    mult2.b <== x3 ;
    x4 <== mult2.c ;
}

component main = three_fac() ;
\end{lstlisting}
Чтобы скомпилировать эту реализацию на Circom в представление R1CS и генератор утверждения, код должен быть сохранён в текстовый файл. Мы будем называть его \fname{three\_fac.circom}. Этот файл затем может быть скомпилирован с помощью следующей команды:
\begin{lstlisting}
circom three_fac.circom --r1cs --wasm --sym
\end{lstlisting}
После компиляции файла команда \texttt{info} может использоваться для вывода статистики схемы, включая количество ограничений (constraints):
\begin{lstlisting}
snarkjs r1cs info circuit.r1cs

[INFO]  snarkJS: Curve: bn-128
[INFO]  snarkJS: # of Wires: 6
[INFO]  snarkJS: # of Constraints: 2
[INFO]  snarkJS: # of Private Inputs: 0
[INFO]  snarkJS: # of Public Inputs: 3
[INFO]  snarkJS: # of Labels: 11
[INFO]  snarkJS: # of Outputs: 1
\end{lstlisting}
\end{example}

\section{Общие концепции программирования}
В этом разделе мы рассмотрим концепции, которые встречаются почти в каждом языке программирования, и увидим, как они могут быть реализованы в компиляторах схем (circuit compilers).
\subsection{Примитивные типы}
% https://zeroknowledge.fm/172-2/ reference for all the languages
Примитивные типы данных, такие как булевы значения (booleans), (беззнаковые) целые числа (unsigned integers) или строки, являются самыми базовыми строительными блоками, которые можно ожидать в каждом универсальном языке программирования высокого уровня. Чтобы писать утверждения в виде компьютерных программ, компилируемых в схемы, необходимо реализовать примитивные типы как системы ограничений (constraint systems) и определить связанные с ними операции как схемы.

В этом разделе мы рассмотрим некоторые распространённые способы достижения этого. После обзора атомарного типа для базового поля, над которым определена схема, мы начнём с реализации булева типа (boolean type) и связанной с ним булевой алгебры как схем. После этого мы определим беззнаковые целые числа (unsigned integers) на основе булева типа и оставим реализацию знаковых целых чисел как упражнение (exercise) для читателя.

\subsubsection{Базовый полевой тип}
\label{def:base_field_type}
Поскольку и алгебраические схемы (algebraic circuits), и связанные с ними системы ограничений ранга~1 определены над конечным полем, элементы этого поля являются атомарными единицами информации в этих моделях. В этом смысле элементы поля $x\in \F$ для алгебраических схем являются тем же, чем биты для компьютеров.

В \lgname{PAPER} мы пишем \texttt{F} для этого типа и указываем фактический экземпляр полевого типа в фигурных скобках после имени данного утверждения. С этим типом связаны две функции, индуцированные операциями \term{addition} (сложения) и \term{multiplication} (умножения) в поле \texttt{F}. Мы пишем
\begin{equation}
\mathtt{MUL}:\; \mathtt{F} \times \mathtt{F} \to \mathtt{F}\;;\; (x,y) \mapsto \mathtt{MUL}(x,y)
\end{equation}
\begin{equation}
\mathtt{ADD}:\; \mathtt{F} \times \mathtt{F} \to \mathtt{F}\;;\; (x,y) \mapsto \mathtt{ADD}(x,y)
\end{equation}
Компиляторы схем (circuit compilers) должны компилировать эти функции в арифметические вентили (gates), как объясняется в~\ref{algebraic-gates}. Каждая другая функция должна быть выражена через эти две атомарные функции.

Чтобы представить сложение и умножение на языке \lgname{PAPER}, мы определяем следующие две функции:
\begin{lstlisting}
fn MUL(x : F, y : F) -> F{}
\end{lstlisting}  
\begin{lstlisting}
fn ADD(x : F, y : F) -> F{}
\end{lstlisting}
Компилятор затем компилирует каждое вхождение функций $\mathtt{MUL}$ или $\mathtt{ADD}$ в следующие графические представления схем:
\begin{center}
\digraph[scale=0.5]{FNMUL}{
	forcelabels=true;
	center=true;
	splines=ortho;
	nodesep= 2.0;
	n1 [shape=box, label="x"];
	n2 [shape=box, label="y"];
	n3 [label="*"];	
	n4 [shape=box, label="(x*y)"];
	
	n1 -> n3 [xlabel="W_1"] ;
	n2 -> n3 [xlabel="W_2"];
	n3 -> n4 [xlabel="W_3"] ;
	
	n5 [shape=box, label="x"];
	n6 [shape=box, label="y"];
	n7 [label="+"];	
	n8 [shape=box, label="(x+y)"];
	
	n5 -> n7 [xlabel="W_1"] ;
	n6 -> n7 [xlabel="W_2"];
	n7 -> n8 [xlabel="W_3"] ;
}
\end{center}
\begin{example}[Базовые вентили] Чтобы дать интуицию о том, как реальный компилятор (compiler) может преобразовать сложение и умножение в алгебраических выражениях в схему, рассмотрим следующее утверждение на \lgname{PAPER}:
\begin{lstlisting}
statement basic_ops {F:F_13} {
  fn main(in_1 : F, pub in_2 : F) -> (F, F){
  	let out_1 : F ;
	let out_2 : F ;
    out_1 <== MUL(in_1,in_2) ;
    out_2 <== ADD(in_1,in_2) ;
    return (out_1, out_2) ;
  }
}
\end{lstlisting} 
Чтобы скомпилировать это в алгебраическую схему (algebraic circuit), мы начинаем с пустой схемы и вычисляем функцию \texttt{main}, которая является единственной функцией в этом утверждении. Мы рисуем подграф входов, содержащий прямоугольные узлы для каждого аргумента функции, и подграф выходов, содержащий прямоугольные узлы для каждого фактора в возвращаемом значении. Поскольку все эти узлы представляют переменные (variables) типа \texttt{field}, нам не нужно добавлять никаких ограничений типов к схеме.

Мы проверяем корректность каждого выражения в правой части каждого оператора \texttt{<==}, включая проверку типов. В нашем случае каждая переменная имеет тип \texttt{field}, и поэтому типы совпадают с типами функций \texttt{MUL} и \texttt{ADD}, а также с типами левых частей операторов \texttt{<==}.

Мы вычисляем выражения в правой части каждого оператора \texttt{<==} индуктивно, заменяя каждое вхождение функции подграфом, представляющим её связанную схему.

Согласно \lgname{PAPER}, каждое вхождение спецификатора instance \texttt{pub} перезаписывает связанное значение witness по умолчанию. Используя соответствующие метки рёбер, мы получаем:
\begin{center}
\digraph[scale=0.5]{BASICOPS}{
    forcelabels=true;
    //center=true;
    splines=ortho;
    nodesep= 2.0;

    subgraph cluster_input {
        nin1 [shape= box, label="in_1"] ;
        nin2 [shape= box, label="in_2"] ;
        color=lightgray ;
        label="inputs" ;
    }

    subgraph cluster_output {
        nout1 [shape= box, label="out_1"] ;
        nout2 [shape= box, label="out_2"] ;
        color=lightgray ;
        label="outputs" ;
    }

    subgraph cluster_mul {
	    nmul1 [shape=box, label="x"];
	    nmul2 [shape=box, label="y"];
	    nmul3 [label="*"];	
	    nmul4 [shape=box, label="(x*y)"];
	    
	    nmul1 -> nmul3 [xlabel="W_1 "] ;
	    nmul2 -> nmul3 [xlabel="W_2"];
	    nmul3 -> nmul4 [xlabel="W_3 "] ;
        color=lightgray ;
        label="fn MUL" ;
    }
    
    subgraph cluster_add {
	    nadd1 [shape=box, label="x"];
	    nadd2 [shape=box, label="y"];
	    nadd3 [label="+"];	
	    nadd4 [shape=box, label="(x+y)"];
	    
	    nadd1 -> nadd3 [xlabel="W_1"] ;
	    nadd2 -> nadd3 [xlabel="W_2 "];
	    nadd3 -> nadd4 [xlabel="W_3 "] ;
        color=lightgray ;
        label="fn ADD" ;
    }    
    // subgraph connectors
    nin1 -> {nmul1, nadd1} [xlabel="W_1", style=dashed, color=grey] ;  
    nin2 -> {nmul2, nadd2} [xlabel="I_1 ", style=dashed, color=grey] ;
    nmul4 -> nout1 [headlabel="W_3 ", style=dashed, color=grey] ;    
    nadd4 -> nout2 [headlabel="W_4 ", style=dashed, color=grey] ;    
}
\end{center}
Любой реальный компилятор (compiler) может обрабатывать связанный с ним высокоуровневый язык подобным образом, заменяя функции или гаджеты (gadgets) предопределёнными связанными схемами. Этот процесс часто сопровождается различными шагами оптимизации (optimization), которые пытаются максимально уменьшить количество ограничений (constraints).

В \lgname{PAPER} мы оптимизируем эту схему, сворачивая все прямоугольные узлы, напрямую связанные с другими прямоугольными узлами, соблюдая правило, что спецификатор \texttt{pub} переменной перезаписывает любой спецификатор witness. Переиндексируя метки рёбер, мы получаем следующую схему как вывод нашего компилятора ``ручка и бумага'':
\begin{center}
\digraph[scale=0.5]{BASICOPSOPTI}{
    forcelabels=true;
    //center=true;
    splines=ortho;
    nodesep= 2.0;

    n1 [shape= box, label="in_1"] ;
    n2 [shape= box, label="in_2"] ;
    n3 [shape= box, label="out_1"] ;
    n4 [shape= box, label="out_2"] ;
	n5 [label="*"] ;	
	n6 [label="+"] ;    
	
    n1 -> {n5, n6} [xlabel="W_1"] ;  
    n2 -> {n5, n6} [xlabel="I_1 "] ;
    n5 -> n3 [xlabel="W_2 "] ;
    n6 -> n4 [label=" W_3"] ;
}
\end{center}
\end{example}
\begin{example}[$3$-факторизация] Рассмотрим нашу задачу $3$-факторизации из~\examplename{}~\ref{ex:3-factorization} и связанную с ней схему $C_{3.fac\_zk}(\F_{13})$, которую мы предоставили в~\examplename{}~\ref{ex:TJJ-circuit_1}. Чтобы понять процесс индуктивной замены высокоуровневых функций их связанными схемами, мы хотим определить утверждение на \lgname{PAPER}, которое ``мозговым компилятором'' (brain-compile) компилируется в алгебраическую схему (algebraic circuit), эквивалентную $C_{3.fac\_zk}(\F_{13})$:
\begin{lstlisting}
statement 3_fac_zk {F:F_13} {
  fn main(x_1 : F, x_2 : F, x_3 : F) -> F{
	let pub 3_fac_zk : F ;
    f_3.fac_zk <== MUL( MUL( x_1 , x_2 ) , x_3 ) ;
    return 3_fac_zk ;
  }
}
\end{lstlisting} 
Используя \lgname{PAPER}, мы начинаем с пустой схемы, а затем добавляем $3$ входных узла в подграф входов, а также $1$ выходной узел в подграф выходов. Все эти узлы помечены связанными именами переменных. Поскольку все эти узлы представляют переменные (variables) типа \texttt{field}, нам не нужно добавлять никаких ограничений типов к схеме.

Мы проверяем корректность каждого выражения в правой части единственного оператора \texttt{<==}, включая проверку типов.

Мы вычисляем выражения в правой части каждого оператора \texttt{<==} индуктивно. У нас есть две вложенные функции умножения, и мы заменяем их связанными схемами умножения, начиная с самой внешней функции. Мы получаем:
\begin{center}
\digraph[scale=0.5]{PAPER3FUC}{
    forcelabels=true;
    //center=true;
    splines=ortho;
    nodesep= 2.0;

    subgraph cluster_input {
        nin1 [shape= box, label="x_1"] ;
        nin2 [shape= box, label="x_2"] ;
        nin3 [shape= box, label="x_3"] ;
        color=lightgray ;
        label="inputs" ;
    }

    subgraph cluster_output {
        nout1 [shape= box, label="f_3.fac_zk"] ;
        color=lightgray ;
        label="outputs" ;
    }

    subgraph cluster_mul1 {
	    nmul11 [shape=box, label="x"];
	    nmul12 [shape=box, label="y"];
	    nmul13 [label="*"];	
	    nmul14 [shape=box, label="(x*y)"];
	    
	    nmul11 -> nmul13 [xlabel="W_1"] ;
	    nmul12 -> nmul13 [xlabel="W_2"];
	    nmul13 -> nmul14 [xlabel="W_3"] ;
        color=lightgray ;
        label="fn MUL" ;
    }
    
    subgraph cluster_mul2 {
	    nmul21 [shape=box, label="x"];
	    nmul22 [shape=box, label="y"];
	    nmul23 [label="*"];	
	    nmul24 [shape=box, label="(x*y)"];
	    
	    nmul21 -> nmul23 [xlabel="W_1"] ;
	    nmul22 -> nmul23 [xlabel="W_2"];
	    nmul23 -> nmul24 [xlabel="W_3"] ;
        color=lightgray ;
        label="fn MUL" ;
    }    
    // subgraph connectors
    nin1 -> nmul11 [xlabel="W_1", style=dashed, color=grey] ;  
    nin2 -> nmul12 [xlabel="W_2", style=dashed, color=grey] ;
    nin3 -> nmul22 [xlabel="W_3", style=dashed, color=grey] ;    
    nmul14 -> nmul21 [xlabel="W_4", style=dashed, color=grey] ;   
    nmul24 -> nout1 [xlabel="I_1", style=dashed, color=grey] ;   
}
\end{center} 
На заключительном шаге оптимизации (optimization) мы сворачиваем все прямоугольные узлы, напрямую связанные с другими прямоугольными узлами, соблюдая правило, что спецификатор \texttt{public} переменной перезаписывает любой спецификатор \texttt{private}. Переиндексируя метки рёбер, мы получаем следующую схему:
\begin{center}
\digraph[scale=0.5]{PAPER3FUCOPTI}{
	forcelabels=true;
	center=true;
	splines=ortho;
	nodesep= 2.0;
	n1 [shape=box, label="x_1"];
	n2 [shape=box, label="x_2"];
	n3 [shape=box, label="x_3"];
	n4 [shape=box, label="f_3.fac_zk"];
	n5 [label="*"];
	n6 [label="*"];	
	
	n1 -> n5 [xlabel="W_1"] ;
	n2 -> n5 [xlabel="W_2"] ;
	n3 -> n6 [xlabel="W_3"];
	n5 -> n6 [xlabel="W_4"] ;
	n6 -> n4 [xlabel="I_1"];
}
\end{center} 
\end{example}
\paragraph{Система ограничений вычитания} По определению, алгебраические схемы содержат только вентили (gates) сложения и умножения, и следовательно, не существует единого вентиля для вычитания в поле, несмотря на то, что вычитание является встроенной операцией в любом поле.

Языки высокого уровня и связанные с ними компиляторы схем (circuit compilers), следовательно, нуждаются в другом способе работы с вычитанием. Чтобы увидеть, как это может быть достигнуто, вспомним, что вычитание определяется сложением с аддитивной обратной, и что обратная может быть эффективно вычислена умножением на $-1$. Схема для полевого вычитания поэтому задаётся
\begin{center}
\digraph[scale=0.4]{BTSUB}{
	forcelabels=true;
	center=true;
	splines=ortho;
	nodesep= 2.0;
	n1 -> n5 [xlabel="S_1    "];
	n2 -> n4 [xlabel="S_2    "];
	n3 -> n4 ;
	n4 -> n5 ;
	n5 -> n6 [xlabel="S_3    "];
	n1 [shape=box, label="x"];
	n2 [shape=box, label="y"];
	n3 [shape=box, label="-1"];
	n4 [label="*"];	
	n5 [label="+"];
	n6 [shape=box, label="SUB(x,y)"]
}
\end{center}
Используя общий метод из~\ref{sec:circuits_associated_R1CS}, связанная с схемой система ограничений ранга~1 задаётся:
\begin{equation}
\left(S_1 + (-1)\cdot S_2\right)\cdot 1 = S_3
\end{equation}
Любое корректное присваивание (assignment) $<S_1,S_2, S_3>$ к этой схеме поэтому принуждает значение $S_3$ быть разностью $S_1- S_2$.

Реальные компиляторы (compilers) обычно предоставляют гаджет (gadget) или функцию для абстрагирования над этой схемой, так что программисты могут использовать вычитание так, как будто оно встроено в схемы.
В \lgname{PAPER} мы определяем следующую функцию вычитания, которая компилируется в предыдущую схему:
\begin{lstlisting}
fn SUB(x : F, y : F) -> F{
  let rslt : F ;
  constant c : F = -1 ;
  rslt <== ADD(x , MUL( y ,  c) );
  return rslt ;
}
\end{lstlisting}
В фазе настройки (setup phase) утверждения мы компилируем каждое вхождение функции $\mathtt{SUB}$ в экземпляр её связанной схемы вычитания, а метки рёбер генерируются согласно правилам из~\ref{def:algebraic-circuit}.
\paragraph{Система ограничений обращения} По определению, алгебраические схемы содержат только вентили (gates) сложения и умножения, и следовательно, не существует единого вентиля для обращения в поле, несмотря на то, что обращение является встроенной операцией в любом поле.

Если лежащее в основе поле является простым полем, один подход состоял бы в использовании малой теоремы Ферма (Fermat's little theorem)~\ref{fermats-little-theorem} для вычисления мультипликативной обратной внутри схемы. Чтобы увидеть, как это работает, пусть $\F_p$ будет простым полем. Мультипликативная обратная $x^{-1}$ элемента поля $x\in\F$ с $x\neq 0$ тогда задаётся как $x^{-1}= x^{p-2}$, и вычисление $x^{p-2}$ в схеме поэтому вычисляет мультипликативную обратную.

К сожалению, реальные простые числа $p$ велики, и вычисление $x^{p-2}$ путём повторного умножения $x$ на само себя неосуществимо. Подход ``квадрирование и умножение'' (square and multiply)~\ref{alg_square-and-mul} быстрее, так как вычисляет степень примерно за $log_2(p)$ шагов, но всё равно добавляет много ограничений (constraints) в схему.

Вычисление обратных в схеме не использует тот факт, что обращение является встроенной операцией в любом поле. Более дружественный к ограничениям (constraints) подход, следовательно, состоит в том, чтобы вычислить мультипликативную обратную вне схемы, а затем только обеспечить корректность вычисления в схеме.

Чтобы понять, как это может быть достигнуто, заметим, что элемент поля $y\in \F$ является мультипликативной обратной элемента поля $x\in \F$ тогда и только тогда, когда $x\cdot y =1$ в $\F$. Мы можем использовать это и определить схему, которая имеет два входа, $x$ и $y$, и принуждает $x\cdot y =1$. Тогда гарантируется, что $y$ является мультипликативной обратной $x$. Цена, которую мы платим, заключается в том, что мы не можем вычислить $y$ путём выполнения схемы, но необходимы вспомогательные данные, чтобы сообщить доказывающему (prover), какое значение $y$ требуется для корректного присваивания (assignment) схемы. Следующая схема определяет ограничение
\begin{center}
\label{circ:inversion}
\digraph[scale=0.4]{BTINV}{
	forcelabels=true;
	center=true;
	splines=ortho;
	nodesep= 2.0;
	n1 -> n3 [xlabel="S_1  "];
	n2 -> n3 [xlabel="S_2  "];
	n3 -> n4 [xlabel="S_3 =1  "];
	n1 [shape=box, label="x"];
	n2 [shape=box, label="x^(-1)"];
	n3 [label="*"];	
	n4 [shape=box, label="1"];	
}
\end{center}
Используя общий метод из~\ref{sec:circuits_associated_R1CS}, схема преобразуется в следующую систему ограничений ранга~1:
\begin{equation}
S_1 \cdot S_2 = 1
\end{equation}
Любое корректное присваивание (assignment) $<S_1,S_2>$ к этой схеме принуждает $S_2$ быть мультипликативной обратной $S_1$, и, поскольку не существует элемента поля $S_2$ такого, что $0\cdot S_2=1$, она также обрабатывает тот факт, что мультипликативная обратная $0$ не определена в любом поле.

Реальные компиляторы (compilers) обычно предоставляют гаджет (gadget) или функцию для абстрагирования над этой схемой, и эти функции вычисляют обратную $x^{-1}$ как часть своего процесса генерации witness. Программистам тогда не нужно беспокоиться о предоставлении обратной как вспомогательных данных схеме. В \lgname{PAPER} мы определяем следующую функцию обращения, которая компилируется в предыдущую схему:
\begin{lstlisting}
fn INV(x : F, y : F) -> F {
  let x_inv : F ;
  constant c : F = 1 ;
  c <== MUL( x ,  y ) ) ;
  x_inv <== y ;
  return x_inv ;
}
\end{lstlisting}
Как мы видим, эта функция принимает два входа: значение поля и его обратную. Следовательно, она не обрабатывает вычисление обратной самостоятельно. Это сделано для того, чтобы сохранить \lgname{PAPER} максимально простым.

В фазе настройки (setup phase) мы компилируем каждое вхождение функции $\mathtt{INV}$ в экземпляр схемы обращения~\ref{circ:inversion}, а метки рёбер генерируются согласно правилам из~\ref{def:algebraic-circuit}.
\paragraph{Система ограничений деления}
\label{def:division_circuit} По определению, алгебраические схемы содержат только вентили (gates) сложения и умножения, и следовательно, не существует единого вентиля для полевого деления, несмотря на то, что деление является встроенной операцией в любом поле.

Реализуя деление как схему, мы используем тот факт, что деление является умножением на мультипликативную обратную. Мы поэтому определяем деление как схему, используя схему обращения и систему ограничений из предыдущего параграфа. Дорогостоящее обращение вычисляется вне схемы, а затем предоставляется как вход схемы. Мы получаем
\begin{center}
\digraph[scale=0.4]{BTDIV}{
	forcelabels=true;
	center=true;
	splines=ortho;
	nodesep= 2.0;
	n1 -> n6 [xlabel="S_1  "];
	n2 -> n4 [xlabel="S_2  "];
	n3 -> n6 [xlabel="S_3  "];
	n3 -> n4 [xlabel="S_3  "];
	n4 -> n5 [xlabel="1  "];
	n6 -> n7 [xlabel="S_4  "];
	n1 [shape=box, label="x"];
	n2 [shape=box, label="y"];
	n3 [shape=box, label="y^(-1)"];
	n4 [label="*"];	
	n5 [shape=box, label="1"];
	n6 [label="*"];	
	n7 [shape=box, label="DIV(x,y)"];	
}
\end{center} 
Используя метод из~\ref{sec:circuits_associated_R1CS}, мы преобразуем эту схему в следующую систему ограничений ранга~1:
\begin{align*}
S_2 \cdot S_3 &= 1\\
S_1 \cdot S_3 &= S_4
\end{align*}
Любое корректное присваивание (assignment) $<S_1,S_2,S_3,S_4>$ к этой схеме принуждает $S_4$ быть полевым делением $S_1$ на $S_2$. Она обрабатывает тот факт, что деление на $0$ не определено, поскольку не существует корректного присваивания в случае $S_2=0$.

В \lgname{PAPER} мы определяем следующую функцию деления, которая компилируется в предыдущую схему:
\begin{lstlisting}
fn DIV(x : F, y : F, y_inv : F) -> F {
  let DIV : F ;
  DIV <== MUL( x ,  INV( y, y_inv ) ) );
  return DIV
}
\end{lstlisting}
В фазе настройки (setup phase) мы компилируем каждое вхождение бинарного оператора $\mathtt{DIV}$ в экземпляр схемы обращения.
\begin{exercise} Пусть $\mathtt{F}$ будет полем $\F_5$ модульной арифметики по модулю $5$ из~\examplename{}~\ref{primfield_z_5}. ``Мозговым компилятором'' (brain-compile) скомпилируйте следующее утверждение на \lgname{PAPER} в алгебраическую схему (algebraic circuit):
\begin{lstlisting}
statement STUPID_CIRC {F: F_5} {
  fn foo(in_1 : F, in_2 : F)->(out_1 : F, out_2 : F,){
    constant c_1 : F = 3 ;
    out_1<== ADD( MUL( c_1 , in_1 ) , in_1 ) ;
    out_2<== INV( c_1 , in_2 ) ;
  } ;
    
  fn main(in_1 : F, in_2 ; F)->(out_1 : F, out_2 : TYPE_2){
	constant (c_1,c_2) : (F,F) = (3,2) ;
    (out_1,out_2) <== foo(in_1, in_2) ;
  } ;
}
\end{lstlisting}
\end{exercise}
\begin{exercise} Рассмотрим кривую tiny-jubjub из~\examplename{}~\ref{TJJ13} и связанную с ней схему~\ref{ex:3-fac-zk-circuit}. Напишите утверждение на \lgname{PAPER}, которое ``мозговым компилятором'' (brain-compile) компилируется в схему, эквивалентную выведенной в~\ref{ex:3-fac-zk-circuit}, предполагая, что точка кривой является instance, а всё остальное присваивание (assignment) --- witness.
\end{exercise}
\begin{exercise} Пусть $\mathtt{F}=\F_{13}$ будет простым полем модульной арифметики по модулю $13$ и $x\in\mathtt{F}$ --- некоторый элемент поля. Определите утверждение на \lgname{PAPER} так, чтобы данный instance $x$ элемент поля $y\in\mathtt{F}$ был witness для утверждения тогда и только тогда, когда $y$ является квадратным корнем из $x$.

``Мозговым компилятором'' (brain-compile) скомпилируйте утверждение в схему и выведите связанную с ней систему ограничений ранга~1. Рассмотрите instance $x=9$ и вычислите конструктивное доказательство для утверждения.
\begin{comment}
\begin{lstlisting}
statement KNOWLEDGE_OF_SQUARE_ROOT {F} {
  fn SQUARE(x : F)->(xx : F){
    xx<== MUL( x , x ) ;
  } ;
    
  fn main(pub x : F, y : F )->(){
    constant c_1 : F = 0 ;
    c_1 <==  SUB( x , SQUARE( y , y ) );
  } ;
}
\end{lstlisting}
\end{comment}
\end{exercise}
\begin{comment}
\paragraph{Modularity} Implementing bounded computation in algebraic circuits it is often necessary to deal with complex expressions of the field type. As we have seen in XXX\sme{add reference} and XXX\sme{add reference}, both algebraic circuits and R1CS have a modularity property, which enables a compiler to derive algebraic circuit implementations for arbitrary circuits.
\begin{example}
To given an intuition of how real-world circuit compilers make use of the modularity property to synthesize complex circuits from high-level programs, we derive a circuit for the following \lgname{PAPER} code:
\begin{lstlisting}
def main<F_13>(private x : F, private y : F, private y^(-1) : F ) -> () {

	circuit:
		outc1 == DIV( ADD( , ) , , );
}
\end{lstlisting}
\end{example}
\begin{example} Consider the prime field $\F_{13}$. In this example, we want to derive an algebraic circuit and associated R1CS that enforces a pair $(x,y)\in \F_{13}^2$ to be the sum of two tiny-jubjub curve points $(x_1,y_1)$ and $(x_2,y_2)$. We assume that we already know that $(x_1,x_2)$ as well as $(x_2,y_2)$ are tiny-jubjub points, that is, we assume that they are the inputs to valid assignments of circuit XXX\sme{add reference}.

To synthesize the associated circuit, we start with the twisted Edwards addition law XXX\sme{add reference} of the tiny-jubjub curve:
$$
(x,y) = \left(\frac{x_1y_2+y_1x_2}{1+8x_1y_1x_2y_2}, \frac{y_1y_2-3x_1x_2}{1-8x_1y_1x_2y_2} \right)
$$ 
To transform this expression into a circuit, we rewrite it in terms of the binary operators $ADD$, $SUB$, $MUL$, $DIV$ that represent the four fundamental field operations in $\F_{13}$. We get
\begin{align*}
(x,y) & = (\\
  & \scriptstyle DIV(ADD(MUL(x_1,y_2),MUL(y_1,x_2)),
         ADD(1,MUL(8,MUL(MUL(x_1,y_1),MUL(x_2,y_2))))), \\   
  & \scriptstyle DIV(ADD(MUL(y_1,y_2),MUL(MUL(3,x_1),x_2)),
         ADD(1,MUL(8,MUL(MUL(x_1,y_1),MUL(x_2,y_2)))))\\
  & )
\end{align*}
We then proceed inductively choosing circuits for the outer most operators, which in this case are two division circuits. We don't expand their inputs into circuits yet, but only represent the inputs symbolically. For better readability we use the symbols of the next operator only, because otherwise the circuit becomes unreadable. We get:
\begin{center}
\digraph[scale=0.4]{TEA}{
	forcelabels=true;
	center=true;
	splines=ortho;
	nodesep= 2.0;
	// x-value
	nx1 -> nx6 [xlabel="Ex_1  "];
	nx2 -> nx4 [xlabel="Ex_2  "];
	nx3 -> nx6 [xlabel="Ex_3  "];
	nx3 -> nx4 [xlabel="Ex_3  "];
	nx4 -> nx5 [xlabel="Ex_4  "];
	nx6 -> nx7 [xlabel="Ex_5  "];
	nx1 [shape=box, label="ADD(.,.)"];
	nx2 [shape=box, label="ADD(.,.)"];
	nx3 [shape=box, label="ADD(.,.)_INV"];
	nx4 [label="*"];	
	nx5 [shape=box, label="1"];
	nx6 [label="*"];	
	nx7 [shape=box, label="DIV(ADD(.,.),ADD(.,.))"];
	// y-value
	ny1 -> ny6 [xlabel="Ey_1  "];
	ny2 -> ny4 [xlabel="Ey_2  "];
	ny3 -> ny6 [xlabel="Ey_3  "];
	ny3 -> ny4 [xlabel="Ey_3  "];
	ny4 -> ny5 [xlabel="Ey_4  "];
	ny6 -> ny7 [xlabel="Ey_5  "];
	ny1 [shape=box, label="ADD(.,.)"];
	ny2 [shape=box, label="ADD(.,.)"];
	ny3 [shape=box, label="ADD(.,.)_INV"];
	ny4 [label="*"];	
	ny5 [shape=box, label="1"];
	ny6 [label="*"];	
	ny7 [shape=box, label="DIV(ADD(.,.),ADD(.,.))"];	
}
\end{center}

\end{example}
\end{comment}
\sme{Is the comment group hidden on purpose?}

\subsubsection{Булев тип}
% implementations can be found here: https://github.com/filecoin-project/zexe/tree/master/snark-gadgets/src/bits
Булевы значения (booleans) --- это классический примитивный тип, реализуемый практически каждым языком программирования более высокого уровня. Поэтому важно реализовать булевы значения в схемах. Одним из самых распространённых способов сделать это является интерпретация аддитивного и мультипликативного нейтральных элементов $\{0,1\}\subset \F$ как двух булевых значений так, что $0$ представляет $false$, а $1$ представляет $true$. Булевы операторы, такие как $and$, $or$ или $xor$, тогда выражаются как алгебраические схемы над $\F$.

Такое представление булевых значений удобно, потому что элементы $0$ и $1$ определены в любом поле. Представление поэтому не зависит от конкретного рассматриваемого поля.

Чтобы зафиксировать нотацию булевой алгебры, мы пишем $0$ для представления $false$ и $1$ для представления $true$, а $\wedge$ пишем для представления булева И (AND), а также $\vee$ для представления булева ИЛИ (OR). Булев оператор НЕ (NOT) записывается как $\lnot$.
\paragraph{Булева система ограничений}
\label{def:boolean_constraint} Чтобы представить булевы значения аддитивным и мультипликативным нейтральными элементами поля, требуется ограничение (constraint), фактически принуждающее переменные булева типа быть либо $1$, либо $0$. На самом деле, многие из следующих схем, представляющих булевы функции, корректны только при условии, что их входные переменные ограничены быть либо $0$, либо $1$. Неограниченность булевых переменных --- это распространённая проблема в проектировании схем.

Чтобы ограничить произвольный элемент поля $x\in \F$ быть $1$ или $0$, ключевое наблюдение состоит в том, что уравнение $x \cdot (1-x) =0$ имеет только два решения $0$ и $1$ в любом поле. Реализация этого уравнения как схемы поэтому генерирует корректное ограничение:
\begin{center}
\digraph[scale=0.4]{BOOLCONS}{
  forcelabels=true;
  center=true;
  splines=ortho;
  nodesep= 2.0;
  nCONS1 -> nCONS4 [xlabel="S_1"] ;
  nCONS1 -> nCONS6 [xlabel="S_1  "] ;
  nCONS2 -> nCONS5 ;
  nCONS3 -> nCONS4 ;
  nCONS4 -> nCONS5 ;
  nCONS5 -> nCONS6 ;
  nCONS6 -> nCONS7 [xlabel="0  "] ;
  nCONS1 [shape=box, label="x"] ;
  nCONS2 [shape=box, label="1"] ;
  nCONS3 [shape=box, label="-1"] ;
  nCONS4 [label="*"] ;
  nCONS5 [label="+"] ;
  nCONS6 [label="*"] ;
  nCONS7 [shape=box, label="0"] ;
}
\end{center}
Используя метод из~\ref{sec:circuits_associated_R1CS}, мы преобразуем эту схему в следующую систему ограничений ранга~1:
$$
S_1 \cdot (1-S_1) = 0
$$
Любое корректное присваивание (assignment) $<S_1>$ к метке этой схемы поэтому принуждает $S_1$ быть либо элементом поля $0$, либо $1$.

Некоторые реальные компиляторы схем (circuit compilers) (такие как \texttt{ZOKRATES} или \texttt{BELLMAN}) являются типизированными, в то время как другие (такие как \texttt{circom}) --- нет. Однако у всех них есть свой способ работы с двоичным ограничением. В \lgname{PAPER} мы определяем булев тип его type\_function, которая компилируется в предыдущую схему:
\begin{lstlisting}
type_function BOOL(b : BOOL) -> F { 
	let x : F ;
	let const c1 : F = 0 ;
	let const c2 : F = 1 ;
	let const c3 : F = -1 ;
    c1 <== MUL( b ,  ADD( c2 , MUL( b , c3) ) ) );
    x <== b ;
    return x ;
}
\end{lstlisting}
В фазе настройки (setup phase) утверждения мы компилируем каждое вхождение переменной булева типа в экземпляр её связанной булевой схемы.
\paragraph{Система ограничений оператора AND} Даны два элемента поля $b_1$ и $b_2$ из $\F$, которые ограничены представлять булевы переменные, мы хотим найти схему, которая вычисляет логический оператор \term{and} (И) $AND(b_1,b_2)$, а также связанную с ней R1CS, принуждающую $b_1$, $b_2$, $AND(b_1,b_2)$ удовлетворять системе ограничений тогда и только тогда, когда $b_1\; \wedge \; b_2 =AND(b_1,b_2)$ выполняется.

Ключевое наблюдение здесь состоит в том, что, даны три булевых ограниченных переменных $b_1$, $b_2$ и $b_3$, уравнение $b_1\cdot b_2 = b_3$ удовлетворяется в $\F$ тогда и только тогда, когда уравнение $b_1\; \wedge \; b_2 = b_3$ удовлетворяется в булевой алгебре. Логический оператор $\wedge$ поэтому реализуем в $\F$ путём полевого умножения его аргументов, и следующая схема вычисляет оператор $\wedge$ в $\F$, предполагая, что все входы ограничены быть $0$ или $1$:
\begin{center}
\digraph[scale=0.4]{BOOLAND}{
  forcelabels=true;
  center=true;
  splines=ortho;
  nodesep= 2.0;
  nAND1 -> nAND3 [xlabel="S_1  "] ;
  nAND2 -> nAND3 [xlabel="S_2"] ;
  nAND3 -> nAND4 [xlabel="S_3  "] ;

  nAND1 [shape=box, label="b_1"] ;
  nAND2 [shape=box, label="b_2"] ;
  nAND3 [label="*"] ;
  nAND4 [shape=box, label="AND(b_1,b_2)"] ;
}
\end{center}
Связанная система ограничений ранга~1 может быть выведена из общего процесса~\ref{sec:circuits_associated_R1CS} и состоит из следующего ограничения:
\begin{equation}
 S_1 \cdot S_2 = S_3
\end{equation}
Распространённые языки схем обычно предоставляют гаджет (gadget) или функцию для абстрагирования над этой схемой, так что программисты могут использовать оператор $\wedge$, не заботясь о связанной схеме. В \lgname{PAPER} мы определяем следующую функцию, которая компилируется в схему оператора $\wedge$:
\begin{lstlisting}
fn AND(b_1 : BOOL, b_2 : BOOL) -> BOOL {
  let AND : BOOL ;
  AND <== MUL( b_1 ,  b_2) ;
  return AND ;
}
\end{lstlisting}
В фазе настройки (setup phase) утверждения мы компилируем каждое вхождение функции $\mathtt{AND}$ в экземпляр связанной с ней схемы оператора $\wedge$.
\paragraph{Система ограничений оператора OR}
\label{def:boolean-or} Даны два элемента поля $b_1$ и $b_2$ из $\F$, которые ограничены представлять булевы переменные, мы хотим найти схему, которая вычисляет логический оператор \term{or} (ИЛИ) $OR(b_1,b_2)$, а также связанную с ней R1CS, принуждающую $b_1$, $b_2$, $OR(b_1,b_2)$ удовлетворять системе ограничений тогда и только тогда, когда $b_1\; \vee \; b_2 = OR(b_1,b_2)$ выполняется.

Предполагая, что три переменные $b_1$, $b_2$ и $b_3$ булево ограничены, уравнение $b_1 + b_2 - b_1\cdot b_2 = b_3$ удовлетворяется в $\F$ тогда и только тогда, когда уравнение $b_1 \; \vee \; b_2 = b_3$ удовлетворяется в булевой алгебре. Логический оператор $\vee$ поэтому реализуем в $\F$ следующей схемой, предполагая, что все входы ограничены быть $0$ или $1$:
\begin{center}
\digraph[scale=0.4]{BOOLOR}{
  forcelabels=true;
  center=true;
  splines=ortho;
  nodesep= 2.0;
  
  nOR1 [shape=box, label="b_1"] ;
  nOR2 [shape=box, label="b_2"] ;
  nOR3 [shape=box, label="-1"] ;
  nOR4 [label="*"] ;
  nOR5 [label="*"] ;
  nOR6 [label="+"] ;
  nOR7 [label="+"] ;
  nOR8 [shape="box", label="OR(b_1,b_2)"]
  
  nOR1 -> nOR4 [xlabel="S_1"] ;
  nOR1 -> nOR6 [xlabel="S_1"] ;
  nOR2 -> nOR4 [taillabel="S_2 "] ;
  nOR2 -> nOR6 [taillabel="S_2 "] ;
  nOR4 -> nOR5 [headlabel="  S_3"] ;
  nOR3 -> nOR5 ; 
  nOR5 -> nOR7 ;
  nOR6 -> nOR7 ;
  nOR7 -> nOR8 [xlabel="S_4 "] ;

}
\end{center}
Связанная система ограничений ранга~1 может быть выведена из общего процесса~\ref{sec:circuits_associated_R1CS} и состоит из следующих ограничений:

\begin{equation}\label{def:boolean-or_constraints}
\begin{tabular}{lcr}
$S_1 \cdot S_2$ & $=$ & $S_3$\\
$(S_1 + S_2 - S_3)\cdot 1$ &$=$ &$S_4$
\end{tabular}
\end{equation}

Распространённые языки схем обычно предоставляют гаджет (gadget) или функцию для абстрагирования над этой схемой, так что программисты могут использовать оператор $\vee$, не заботясь о связанной схеме. В \lgname{PAPER} мы определяем следующую функцию, которая компилируется в схему оператора $\vee$:
\begin{lstlisting}
fn OR(b_1 : BOOL, b_2 : BOOL) -> BOOL {
  let OR : BOOL ;
  let const c1 : F = -1 ;
  OR <== ADD(ADD(b_1,b_2),MUL(c1,MUL(b_1,b_2))) ;
  return OR ;
}
\end{lstlisting}
В фазе настройки (setup phase) утверждения мы компилируем каждое вхождение функции $\mathtt{OR}$ в экземпляр связанной с ней схемы оператора $\vee$.
\begin{exercise} Пусть $\F$ будет конечным полем и пусть $b_1$, а также $b_2$ --- две булево ограниченные переменные из $\F$. Покажите, что уравнение
$OR(b_1,b_2) = 1 - (1 - b_1)\cdot (1 - b_2)$ выполняется.

Используйте это уравнение для выведения алгебраической схемы с входящими переменными $b_1$ и $b_2$ и исходящей переменной $OR(b_1,b_2)$ так, чтобы $b_1$ и $b_2$ были булево ограничены и схема имела корректное присваивание (assignment) тогда и только тогда, когда $OR(b_1,b_2) = b_1 \vee b_2$.

Используйте технику из~\ref{sec:circuits_associated_R1CS} для преобразования этой схемы в систему ограничений ранга~1 и найдите её полное множество решений. Определите функцию на \lgname{PAPER}, которая ``мозговым компилятором'' (brain-compile) компилируется в схему.
\begin{comment}
\begin{center}
\digraph[scale=0.4]{BOOLOR}{
  forcelabels=true;
  center=true;
  splines=ortho;
  nodesep= 2.0;
  nOR1 -> nOR5 [xlabel="S_1  "] ;
  nOR2 -> nOR7 [xlabel="S_2  "] ;
  nOR3 -> {nOR5, nOR7, nOR10} ;
  nOR4 -> {nOR6, nOR8, nOR11} ;
  nOR5 -> nOR6; 
  nOR6 -> nOR9 ;
  nOR7 -> nOR8 ;
  nOR8 -> nOR9 ;
  nOR9 -> nOR10 [xlabel="S_3  "] ;
  nOR10 -> nOR11 ;
  nOR11 -> nOR12 [xlabel="S_4  "] ;

  nOR1 [shape=box, label="b_1"] ;
  nOR2 [shape=box, label="b_2"] ;
  nOR3 [shape=box, label="-1"] ;
  nOR4 [shape=box, label="1"] ;
  nOR5 [label="*"] ;
  nOR6 [label="+"] ;
  nOR7 [label="*"] ;
  nOR8 [label="+"] ;
  nOR9 [label="*"] ;
  nOR10 [label="*"] ;
  nOR11 [label="+"] ;
  nOR12 [shape=box, label="OR(b_1,b_2)"] ;
}
\end{center}
The associated Rank-1 Constraint System can be deduced from the general process \ref{sec:circuits_associated_R1CS} and consists of the following constraints
\begin{align*}
 (1- S_1) \cdot (1-S_2) & = S_3\\
  (1-S_3)\cdot 1 &= S_4
\end{align*}
Common circuit languages typically provide a gadget or a function to abstract over this circuit such that programers can use the $\vee$ operator without caring about the associated circuit. In \lgname{PAPER}, we define the following function that compiles to the $\vee$-operator's circuit:
\begin{lstlisting}
fn OR(b_1 : BOOL, b_2 : BOOL) -> OR(b_1,b_2) : BOOL{
  l
  constant c1 = 1 ;
  constant c2 = -1 ;
  OR(b_1,b_2) <== ADD(c1,MUL(MUL(ADD(c1,MUL(b_1,c2)),ADD(c1,MUL(b_1,c2))),c2))  ;
}
\end{lstlisting}
In the setup phase of a statement, we compile every occurrence of the $\mathtt{OR}$ function into an instance of its associated $\vee$-operator's circuit.
\end{comment} 
\end{exercise}
\paragraph{Система ограничений оператора NOT} Дан элемент поля $b$ из $\F$, который ограничен представлять булеву переменную, мы хотим найти схему, которая вычисляет логический оператор \term{NOT} (НЕ) $NOT(b)$, а также связанную с ней R1CS, принуждающую $b$, $NOT(b)$ удовлетворять системе ограничений тогда и только тогда, когда $\lnot b = NOT(b)$ выполняется.

Предполагая, что две переменные $b_1$ и $b_2$ булево ограничены, уравнение $(1-b_1) = b_2$ удовлетворяется в $\F$ тогда и только тогда, когда уравнение $\lnot b_1 = b_2$ удовлетворяется в булевой алгебре. Логический оператор $\lnot$ поэтому реализуем в $\F$ следующей схемой, предполагая, что все входы ограничены быть $0$ или $1$:
\begin{center}
\digraph[scale=0.4]{BOOLNOT}{
  forcelabels=true;
  center=true;
  splines=ortho;
  nodesep= 2.0;
  nNOT1 -> nNOT4 [xlabel="S_1  "] ;
  nNOT2 -> nNOT4 ;
  nNOT3 -> nNOT5 ;
  nNOT4 -> nNOT5 ;
  nNOT5 -> nNOT6 [xlabel="S_2  "] ;

  nNOT1 [shape=box, label="b"] ;
  nNOT2 [shape=box, label="-1"] ;
  nNOT3 [shape=box, label="1"] ;
  nNOT4 [label="*"] ;
  nNOT5 [label="+"] ;
  nNOT6 [shape=box, label="NOT(b)"] ;
}
\end{center}
Связанная система ограничений ранга~1 может быть выведена из общего процесса~\ref{sec:circuits_associated_R1CS} и состоит из следующих ограничений
\begin{align*}
  (1-S_1)\cdot 1 &= S_2
\end{align*}
Распространённые языки схем обычно предоставляют гаджет (gadget) или функцию для абстрагирования над этой схемой, так что программисты могут использовать оператор $\lnot$, не заботясь о связанной схеме. В \lgname{PAPER} мы определяем следующую функцию, которая компилируется в схему оператора $\lnot$:
\begin{lstlisting}
fn NOT(b : BOOL) -> BOOL{
  let NOT : BOOL ;
  let const c1 = 1 ;
  let const c2 = -1 ;
  NOT <== ADD( c1 , MUL( c2 , b) ) ;
  return NOT ;
}
\end{lstlisting}
В фазе настройки (setup phase) утверждения мы компилируем каждое вхождение функции $\mathtt{NOT}$ в экземпляр связанной с ней схемы оператора $\lnot$.
\begin{exercise}
Пусть $\F$ будет конечным полем. Выведите алгебраические схемы и связанные с ними системы ограничений ранга~1 для следующих операторов: NOR, XOR, NAND, EQU.
\end{exercise}
\paragraph{Модульность} Как мы видели в главе~\ref{sec:statements}, и алгебраические схемы, и R1CS обладают свойством модульности (modularity), и как мы видели в этом разделе, все базовые булевы функции выражаемы в схемах. Комбинация этих двух свойств показывает, что возможно выражать произвольные булевы функции как алгебраические схемы.

Выразительность алгебраических схем, а следовательно, систем ограничений ранга~1 столь же обща, как выразительность булевых схем. Важное следствие состоит в том, что языки $L_{R1CS-SAT}$ и $L_{Circuit-SAT}$, как определены в~\ref{r1cs-satisfiability} и~\ref{circuit-satisfiability}, столь же общи, как известный язык $L_{3-SAT}$, который известен как $\mathcal{NP}$-полный.
\begin{example} Чтобы привести пример того, как компилятор (compiler) может конструировать сложные булевы выражения в алгебраических схемах из простых, и как мы выводим их связанные системы ограничений ранга~1, рассмотрим следующее утверждение на \lgname{PAPER}:
\begin{lstlisting}
statement BOOLEAN_STAT {F: F_p} {
  fn main(b_1:BOOL,b_2:BOOL,b_3:BOOL,b_4:BOOL )-> BOOL {
    let pub b_5 : BOOL ;
    b_5 <== AND( OR( b_1 , b_2) , AND( b_3 , NOT( b_4) ) ) ;
    return b_5 ;
  } ;
}
\end{lstlisting}
Код описывает схему, которая принимает четыре witness-входные переменные $b_1$, $b_2$, $b_3$ и $b_4$ булева типа и вычисляет instance-переменную выхода $b_5$ так, чтобы следующее булево выражение выполнялось:
$$
\left( b_1 \vee b_2 \right) \wedge (b_3 \wedge \lnot b_4) = b_5
$$
Во время фазы настройки (setup-phase) компилятор схем (circuit compiler) преобразует это высокоуровневое утверждение в схему и связанную систему ограничений ранга~1, тем самым определяя язык $L_{BOOLEAN\_STAT}$.

Чтобы увидеть, как это может быть достигнуто, мы используем \lgname{PAPER} как пример для выполнения фазы настройки (setup-phase) и компиляции \texttt{BOOLEAN\_STAT} в схему. Учитывая определение булева ограничения~\ref{def:boolean_constraint}, а также определения соответствующих булевых операторов, мы получаем следующую схему:
\begin{center}
\digraph[scale=0.25]{BOOLCOMPLEX}{
  forcelabels=true;
  center=true;
  splines=ortho;
  nodesep= 1.0;
  rankdir=TB;

  subgraph clusterINPUT {
    nb1 [shape=box, label="b_1"] ;
    nb2 [shape=box, label="b_2"] ;    
    nb3 [shape=box, label="b_3"] ;
    nb4 [shape=box, label="b_4"] ;
    color=lightgray;
    label="inputs" ; 
  }
  
  subgraph clusterOUTPUT {
    nout1 [shape=box, label="AND( OR( b1 , b2 ) , AND( b3 , NOT( b4 ) )"] 
    color=lightgray;
    label="outputs" ; 
  }
  
  subgraph clusterCONS  {
    subgraph clusterBCONS1 {
      nCONSB11 -> nCONSB14 [xlabel="S_1"] ;
      nCONSB11 -> nCONSB16 [xlabel="S_1"] ;
      nCONSB12 -> nCONSB15 ;
      nCONSB13 -> nCONSB14 ;
      nCONSB14 -> nCONSB15 ;
      nCONSB15 -> nCONSB16 ;
      nCONSB16 -> nCONSB17 [xlabel="w_5=0  "] ;
      nCONSB11 [shape=box, label="b", color=lightgray ] ;
      nCONSB12 [shape=box, label="1"] ;
      nCONSB13 [shape=box, label="-1"] ;
      nCONSB14 [label="*"] ;
      nCONSB15 [label="+"] ;
      nCONSB16 [label="*"] ;
      nCONSB17 [shape=box, label="0"] ;
      color=lightgray;
      label="b1 : BOOL" ;
    }
    
      subgraph clusterBCONS2 {
      nCONSB21 -> nCONSB24 [xlabel="S_1"] ;
      nCONSB21 -> nCONSB26 [xlabel="S_1"] ;
      nCONSB22 -> nCONSB25 ;
      nCONSB23 -> nCONSB24 ;
      nCONSB24 -> nCONSB25 ;
      nCONSB25 -> nCONSB26 ;
      nCONSB26 -> nCONSB27 [xlabel="W_6=0  "] ;
      nCONSB21 [shape=box, label="b", color=lightgray ] ;
      nCONSB22 [shape=box, label="1"] ;
      nCONSB23 [shape=box, label="-1"] ;
      nCONSB24 [label="*"] ;
      nCONSB25 [label="+"] ;
      nCONSB26 [label="*"] ;
      nCONSB27 [shape=box, label="0"] ;
      color=lightgray;
      label="b2 : BOOL" ;
    }
    
    subgraph clusterBCONS3 {
      nCONSB31 -> nCONSB34 [xlabel="S_1"] ;
      nCONSB31 -> nCONSB36 [xlabel="S_1"] ;
      nCONSB32 -> nCONSB35 ;
      nCONSB33 -> nCONSB34 ;
      nCONSB34 -> nCONSB35 ;
      nCONSB35 -> nCONSB36 ;
      nCONSB36 -> nCONSB37 [xlabel="W_7=0  "] ;
      nCONSB31 [shape=box, label="b", color=lightgray ] ;
      nCONSB32 [shape=box, label="1"] ;
      nCONSB33 [shape=box, label="-1"] ;
      nCONSB34 [label="*"] ;
      nCONSB35 [label="+"] ;
      nCONSB36 [label="*"] ;
      nCONSB37 [shape=box, label="0"] ;
      color=lightgray;
      label="b3 : BOOL" ;
    }
    
    subgraph clusterBCONS4 {
      nCONSB41 -> nCONSB44 [xlabel="S_1"] ;
      nCONSB41 -> nCONSB46 [xlabel="S_1"] ;
      nCONSB42 -> nCONSB45 ;
      nCONSB43 -> nCONSB44 ;
      nCONSB44 -> nCONSB45 ;
      nCONSB45 -> nCONSB46 ;
      nCONSB46 -> nCONSB47 [xlabel="W_8=0  "] ;
      nCONSB41 [shape=box, label="b", color=lightgray ] ;
      nCONSB42 [shape=box, label="1"] ;
      nCONSB43 [shape=box, label="-1"] ;
      nCONSB44 [label="*"] ;
      nCONSB45 [label="+"] ;
      nCONSB46 [label="*"] ;
      nCONSB47 [shape=box, label="0"] ;
      color=lightgray;
      label="b4 : BOOL" ;
    } 
    color = white ; 
  } 

  subgraph clusterCIRC{

    subgraph clusterORb1b2 {
      nOR1 [shape=box, label="b_1", color=lightgray] ;
      nOR2 [shape=box, label="b_2", color=lightgray] ;
      nOR3 [shape=box, label="-1"] ;
      nOR4 [label="*"] ;
      nOR5 [label="*"] ;
      nOR6 [label="+"] ;
      nOR7 [label="+"] ;
      nOR8 [shape="box", label="OR(b_1,b_2)", color=lightgray]
      
      nOR1 -> nOR4 [xlabel="S_1"] ;
      nOR1 -> nOR6 [xlabel="S_1"] ;
      nOR2 -> nOR4 [taillabel="S_2 "] ;
      nOR2 -> nOR6 [taillabel="S_2 "] ;
      nOR4 -> nOR5 [headlabel="  S_3"] ;
      nOR3 -> nOR5 ; 
      nOR5 -> nOR7 ;
      nOR6 -> nOR7 ;
      nOR7 -> nOR8 [xlabel="S_4 "] ;
      color=lightgray;
      label="fn OR" ;
    }

    subgraph clusterANDb3NOTb4 {
      nAND21 -> nAND23 [headlabel="  S_1"] ;
      nAND22 -> nAND23 [xlabel="S_2"] ;
      nAND23 -> nAND24 [xlabel="S_3 "] ;

      nAND21 [shape=box, label="b_1", color=lightgray ] ;
      nAND22 [shape=box, label="b_2", color=lightgray ] ;
      nAND23 [label="*"] ;
      nAND24 [shape=box, label="AND( b_1 , b_2 )", color=lightgray] ;
      color=lightgray;
      label="fn AND"
    }

    subgraph clusterNOTb4 {
      nNOT1 -> nNOT4 [xlabel="S_1 "] ;
      nNOT2 -> nNOT4 ;
      nNOT3 -> nNOT5 ;
      nNOT4 -> nNOT5 ;
      nNOT5 -> nNOT6 [headlabel="S_2 "] ;

      nNOT1 [shape=box, label="b", color=lightgray ] ;
      nNOT2 [shape=box, label="-1"] ;
      nNOT3 [shape=box, label="1"] ;
      nNOT4 [label="*"] ;
      nNOT5 [label="+"] ;
      nNOT6 [shape=box, label="NOT( b )", color=lightgray ] ;
      color=lightgray;
      label="fn NOT"
    }

    subgraph clusterAND1 {
      nAND1_1 -> nAND1_3 [xlabel="S_1"] ;
      nAND1_2 -> nAND1_3 [xlabel="S_2 "] ;
      nAND1_3 -> nAND1_4 [xlabel="S_3 "] ;
      nAND1_1 [shape=box, label="b_1", color=lightgray ] ;
      nAND1_2 [shape=box, label="b_2", color=lightgray] ;
      nAND1_3 [label="*"] ;
      nAND1_4 [shape=box, label="AND( b_1 , b_2 )", color=lightgray] ;
      color=lightgray;
      label="fn AND"
    }
    nNOT6 -> nAND22 [style=dashed, color=grey] ;
    nOR8 -> nAND1_1 [style=dashed, color=grey] ;
    nAND24 -> nAND1_2 [style=dashed, color=grey] ;
    
    color = white ; 
  }
  // outer circuit
    nb1 -> nCONSB11 [xlabel="W_1", style=dashed, color=grey] ;
    nCONSB11 -> nOR1 [style=dashed, color=grey] ;
    nb2 -> nCONSB21 [headlabel="W_2", style=dashed, color=grey] ;
    nCONSB21 -> nOR2 [style=dashed, color=grey] ;
    nb3 -> nCONSB31 [xlabel="W_3 ", style=dashed, color=grey] ;
    nCONSB31 -> nAND21 [style=dashed, color=grey] ;
    nb4 -> nCONSB41 [xlabel="W_4", style=dashed, color=grey] ;
    nCONSB41 -> nNOT1 [style=dashed, color=grey] ;
    nAND1_4 -> nout1 [xlabel="I_1"; style=dashed, color=grey] ;

}
\end{center}\sme{can we rotate this by $90^{\circ}$? Good question. IDK}
Простая оптимизация (optimization) затем сворачивает все прямоугольные узлы, напрямую связанные, и все прямоугольные узлы, представляющие одинаковые константы (constants). После перемаркировки рёбер следующая схема представляет схему, связанную с утверждением \texttt{BOOLEAN\_STAT}:
\begin{center}
\digraph[scale=0.5]{BOOLCOMPLEXOPTI}{
  forcelabels=true;
  center=true;
  splines=ortho;

  one -> nCONSb15 ;
  minusone -> nCONSb14 ;
  nCONSb14 -> nCONSb15 ;
  nCONSb15 -> nCONSb16 ;
  nCONSb16 -> zero [xlabel="W_5"] ;
  nCONSb14 [label="*"] ;
  nCONSb15 [label="+"] ;
  nCONSb16 [label="*"] ;

  one -> nCONSb25 ;
  minusone -> nCONSb24 ;
  nCONSb24 -> nCONSb25 ;
  nCONSb25 -> nCONSb26 ;
  nCONSb26 -> zero [headlabel="W_6 "] ;
  nCONSb24 [label="*"] ;
  nCONSb25 [label="+"] ;
  nCONSb26 [label="*"] ;

  one -> nCONSb35 ;
  minusone -> nCONSb34 ;
  nCONSb34 -> nCONSb35 ;
  nCONSb35 -> nCONSb36 ;
  nCONSb36 -> zero [taillabel="W_7"] ;
  nCONSb34 [label="*"] ;
  nCONSb35 [label="+"] ;
  nCONSb36 [label="*"] ;

  one -> nCONSb45 ;
  minusone -> nCONSb44 ;
  nCONSb44 -> nCONSb45 ;
  nCONSb45 -> nCONSb46 ;
  nCONSb46 -> zero [taillabel="W_8 "] ;
  nCONSb44 [label="*"] ;
  nCONSb45 [label="+"] ;
  nCONSb46 [label="*"] ;

  //minusone -> {nOR5, nOR7, nOR10} ;
  //one -> {nOR6, nOR8, nOR11} ;
  //nOR5 -> nOR6; 
  //nOR6 -> nOR9 ;
  //nOR7 -> nOR8 ;
  //nOR8 -> nOR9 ;
  //nOR9 -> nOR10 [headlabel="W_9  "] ;
  //nOR10 -> nOR11 ;
  //nOR5 [label="*"] ;
  //nOR6 [label="+"] ;
  //nOR7 [label="*"] ;
  //nOR8 [label="+"] ;
  //nOR9 [label="*"] ;
  //nOR10 [label="*"] ;
  //nOR11 [label="+"] ;
  
  // new 
  //nOR1 [shape=box, label="b_1"] ;
  //nOR2 [shape=box, label="b_2"] ;
  //nOR3 [shape=box, label="-1"] ;
  minusone -> nOR5 ;
  nOR4 [label="*"] ;
  nOR5 [label="*"] ;
  nOR6 [label="+"] ;
  nOR7 [label="+"] ;
  //nOR8 [shape="box", label="OR(b_1,b_2)"]
  
  //nOR1 -> nOR4 [xlabel="S_1"] ;
  //nOR1 -> nOR6 [xlabel="S_1"] ;
  //nOR2 -> nOR4 [taillabel="S_2 "] ;
  //nOR2 -> nOR6 [taillabel="S_2 "] ;
  nOR4 -> nOR5 [taillabel="W_9"] ;
  //nOR3 -> nOR5 ; 
  nOR5 -> nOR7 ;
  nOR6 -> nOR7 ;
  //nOR7 -> nOR8 [xlabel="S_4 "] ;

  nAND23 [label="*"] ;
  minusone -> nNOT4 ;
  one -> nNOT5 ;
  nNOT4 -> nNOT5 ;
  nNOT4 [label="*"] ;
  nNOT5 [label="+"] ;

  nOR7 -> nAND1_3; // here
  nAND23 -> nAND1_3 [xlabel="W_10"] ;
  nAND1_3 -> nAND1_4 [xlabel="I_1 "] ;
  nAND1_3 [label="*"] ;
  nAND1_4 [shape=box, label="AND( OR( b1 , b2 ) , AND( b3 , NOT( b4 ) )"] ;

  // outer circuit
    nNOT5 -> nAND23 ;
    b1 -> nOR4 [taillabel="  W_1"] ;
    b1 -> nOR6 [xlabel="W_1"] ;
    b1 -> nCONSb14 [headlabel="W_1 "] ;
    b1 -> nCONSb16 [headlabel="W_1     "] ;
    b2 -> nOR4 [taillabel=" W_2"] ;
    b2 -> nOR6  [headlabel=""] ;
    b2 -> nCONSb24 [headlabel=" W_2 "] ;
    b2 -> nCONSb26 [headlabel=" W_2 "] ;
    b3 -> nAND23 [xlabel=" W_3 "] ;
    b3 -> nCONSb34 [taillabel=" W_3"] ;
    b3 -> nCONSb36 [taillabel=""] ;
    b4 -> nNOT4 [headlabel=" W_4"] ;
    b4 -> nCONSb44 [headlabel=" W_4"] ;
    b4 -> nCONSb46 [headlabel="W_4"] ;
    b1 [shape=box, label="b1"] ;
    b2 [shape=box, label="b2"] ;
    b3 [shape=box, label="b3"] ;
    b4 [shape=box, label="b4"] ;
    minusone [shape=box, label="-1"] ;
    one [shape=box, label="1"] ;
    zero [shape=box, label="0"] ;
}
\end{center}
Для некоторой instance-переменной $I_1$ из $\F_{13}$ корректное присваивание (assignment) к этой схеме состоит из witness-переменных $W_1$, $W_2$, $W_3$, $W_4$ из $\F_{13}$ таких, что уравнение $I_1 = \left( W_1 \vee W_2 \right) \wedge (W_3 \wedge \lnot W_4)$ выполняется. Кроме того, корректное присваивание также должно содержать witness-переменные $W_5$, $W_6$, $W_7$, $W_8$, $W_9$ и $W_{10}$, которые могут быть выведены из выполнения схемы. Переменные $W_5$, $\ldots$, $W_8$ обеспечивают, что первые четыре witness-переменные ограничены быть либо $0$, либо $1$, но не каким-либо другим элементом поля, а остальные обеспечивают булевы операции в выражении.

Чтобы вычислить связанную R1CS, мы можем использовать общий метод из~\ref{sec:circuits_associated_R1CS} и посмотреть на каждое помеченное исходящее ребро, не идущее от источникового узла, в оптимизированной схеме. Мы объявляем ребро, идущее к единственному выходному узлу, как instance, а каждое другое ребро --- как witness. В этом случае мы получаем:
\begin{align*}
W_5:\;\; & W_1 \cdot (1- W_1) = 0  & \text{булевы ограничения}\\
W_6:\;\; & W_2 \cdot (1- W_2) = 0 \\
W_7:\;\; & W_3 \cdot (1- W_3) = 0 \\
W_8:\;\; & W_4 \cdot (1- W_4) = 0 \\
W_9:\;\; & W_1 \cdot W_2 = W_9 & \text{ первое ограничение OR-оператора}\\
W_{10}:\;\; & W_3 \cdot (1-W_4) = W_{10} & \text{ограничения AND(.,NOT(.))-оператора}\\
I_1:\;\; & (W_1 + W_2 -W_9) \cdot W_{10} = I_1 & \text{ограничения AND-оператора}\\
\end{align*}
Причина, по которой эта R1CS содержит только одно ограничение для вентиля (gate) умножения в OR-схеме, в то время как общее определение~\ref{def:boolean-or} требует двух ограничений, состоит в том, что второе ограничение в~\ref{def:boolean-or_constraints} появляется только потому, что конечный вентиль сложения связан с выходным узлом. В этом случае, однако, конечный вентиль сложения из OR-схемы обеспечен в левом множителе ограничения $I_{1}$. Аналогичное верно и для схемы отрицания.

Во время фазы доказывающего (prover-phase) должно быть дано некоторое публичное instance $I_1$. Чтобы вычислить конструктивное доказательство для утверждения связанных языков относительно instance $I_1$, доказывающий (prover) должен найти четыре булевых значения $W_1$, $W_2$, $W_3$ и $W_4$ таких, что
$$
\left( W_1 \vee W_2 \right) \wedge (W_3 \wedge \lnot W_4) = I_1
$$ 
выполняется. В нашем случае ни схема, ни утверждение на \lgname{PAPER} не указывают, как найти эти значения, и это проблема, которую любой доказывающий (prover) должен решить вне схемы. Это может быть или не быть верно и для других задач. В любом случае, как только доказывающий (prover) нашёл эти значения, он может выполнить схему, чтобы найти корректное присваивание (assignment).

Чтобы привести конкретный пример, пусть $I_1=1$ и предположим $W_1=1$, $W_2=0$, $W_3=1$ и $W_4=0$. Поскольку
$\left( 1 \vee 0 \right) \wedge (1 \wedge \lnot 0) = 1$, эти значения удовлетворяют задаче, и мы можем использовать их для выполнения схемы. Мы получаем
\begin{align*}
W_5 & = W_1 \cdot (1- W_1) = 0\\
W_6 & = W_2 \cdot (1- W_2) = 0 \\
W_7 & = W_3 \cdot (1- W_3) = 0 \\
W_8 & = W_4 \cdot (1- W_4) = 0 \\
W_9 & = W_1\cdot W_2 = 0\\
W_{10} & = W_3 \cdot (1-W_4) = 1\\
I_1 & = (W_1 + W_2 - W_9) \cdot W_{10} = 1
\end{align*}
Конструктивное доказательство знания witness, например, для $I_1=1$, поэтому задаётся строкой $\pi=<W_5,W_6,W_7,W_8,W_9,W_{10}>=<0,0,0,0,0,1>$.
\end{example}
\subsubsection{Массивы} Тип \texttt{array} (массив) представляет собой коллекцию фиксированного размера элементов одинакового типа, каждый из которых выбирается одним или несколькими индексами, которые могут вычисляться во время выполнения программы.

Массивы --- это классический тип, реализуемый многими языками программирования более высокого уровня, компилируемыми в схемы или системы ограничений ранга~1. Однако большинство высокоуровневых языков схем поддерживают \term{static} (статические) массивы, т.е. массивы, длина которых известна только во время компиляции.

Наиболее распространённый способ компиляции массивов в схемы состоит в преобразовании любого массива данного типа $\texttt{t}$ и размера $\texttt{N}$ в $\texttt{N}$ переменных (variables) схемы типа $\texttt{t}$. Массивы поэтому являются \term{syntactic sugar} (синтаксическим сахаром), цель которого --- сделать код более читаемым и записываемым для человека. В \lgname{PAPER} мы определяем следующую type\_function массива:
\begin{lstlisting}
type_function <Name>: <Type>[N : unsigned] -> (Type,...) { 
  return (<Name>[0],...)
}
\end{lstlisting}
В фазе настройки (setup phase) утверждения мы компилируем каждое вхождение массива размера \texttt{N}, содержащего элементы типа \texttt{Type}, в \texttt{N} переменных (variables) типа \texttt{Type}.
\begin{example} Чтобы дать интуицию о том, как реальный компилятор (compiler) может преобразовать массивы в переменные (variables) схемы, рассмотрим следующее утверждение на \lgname{PAPER}:
\begin{lstlisting}
statement ARRAY_TYPE {F: F_5} {
  fn main(x: F[2])-> F {
  	let constant c: F[2] = [2,4] ;
  	let out : F <== MUL(ADD(x[1],c[0]),ADD(x[0],c[1])) ;
  	return out ;
  } ;
}
\end{lstlisting}
Во время фазы настройки (setup phase) компилятор схем (circuit compiler) может затем заменить любое вхождение типа массива кортежем переменных (variables) лежащего в основе типа, а затем использовать эти переменные в процессе синтеза схемы. Чтобы увидеть, как это может быть достигнуто, мы используем \lgname{PAPER} как пример. Абстрагируясь над подсхемой вычисления, мы получаем следующую схему:
\begin{center}
\digraph[scale=0.5]{SIMPLEARRAYCIRC}{
  forcelabels=true;
  center=true;
  splines=ortho;
  nodesep= 2.0;

    subgraph cluster_input {
      subgraph cluster_array1 {
          nx1 [shape= box, label="x[0]"] ;
          nx2 [shape= box, label="x[1]"] ;
          color=lightgray ;
          label="x[]" ;
      }

      subgraph cluster_array2 {
          nc1 [shape= box, label="c[0]"] ;
          nc2 [shape= box, label="c[1]"] ;
          color=lightgray ;
          label="c[]" ;
      }       
      color=lightgray ;
      label="inputs" ;
    }

    subgraph cluster_output {
        nout1 [shape= box, label="out"] ;
        color=lightgray ;
        label="outputs" ;
    }

    subgraph cluster_expression {
	    nexp1 [shape=box, label="MUL(ADD,ADD)"];
      color=lightgray ;
      label="fn MUL" ;
    }
   
    // subgraph connectors
    nx1 -> nexp1 [xlabel="Wx0"] ;  
    nx2 -> nexp1 [xlabel="Wx1"] ; 
    nc1 -> nexp1 [xlabel="Wc0"] ; 
    nc2 -> nexp1 [xlabel="Wc1"] ; 
    nexp1 -> nout1 [xlabel="Wout"] ;    
}
\end{center}
\end{example}
\subsubsection{Беззнаковый целочисленный тип} Беззнаковые целые числа (unsigned integers) размера \texttt{N}, где \texttt{N} обычно является степенью двойки, представляют неотрицательные целые числа в диапазоне $0\ldots 2^N-1$. У них есть понятие сложения, вычитания и умножения, определённое модульной арифметикой по модулю $2^N$. Если дано некоторое \texttt{N}, мы пишем \texttt{uN} для связанного типа.

\paragraph{Система ограничений uN} Многие высокоуровневые языки схем определяют различные типы \texttt{uN} как массивы размера \texttt{N}, где каждый элемент имеет булев тип. Это параллельно их представлению на обычном компьютерном оборудовании и позволяет эффективно и прямолинейно определять распространённые операторы, такие как различные операторы сдвига или логические операторы.

Предполагая, что некоторое беззнаковое целое число (unsigned integer) \texttt{N} известно во время компиляции в \lgname{PAPER}, мы определяем следующую type\_function \texttt{uN}, которая приводит беззнаковое целое число (unsigned integer) к массиву булевых переменных:
\begin{lstlisting}
type_function uN -> BOOL[N] { 
  let base2 : BOOL[N] ;
  base2[0] <== uN[0] ;
  base2[1] <== uN[1] ;
  ...
  base2[N] <== uN[N] ; 
  return base2 ;
}
\end{lstlisting}
Чтобы принудить $N$-куплет элементов поля $<b_0,\ldots,b_{N-1}>$ представлять элемент типа \texttt{uN}, нам поэтому требуются $N$ булевых ограничений
\begin{align*}
S_0 \cdot (1-S_0) & = 0\\
S_1 \cdot (1-S_1) & = 0\\
\cdots &\\
S_{N-1} \cdot (1-S_{N-1}) & = 0\\
\end{align*}
\begin{remark} Представление типа \texttt{uN} как булевых массивов концептуально чисто и работает над произвольными базовыми полями. Однако представление беззнаковых целых чисел (unsigned integers) таким образом требует много места, поскольку каждый бит представляется как элемент поля, и если базовое поле велико, эти элементы поля требуют значительного места в оборудовании.

Следует отметить, что в некоторых случаях существует другой, более эффективный с точки зрения места и системы ограничений (Constraint System) подход для представления беззнаковых целых чисел (unsigned integers), который может использоваться всякий раз, когда лежащее в основе базовое поле достаточно велико. Чтобы понять это, вспомним, что сложение и умножение в простом поле $\F_p$ равны сложению и умножению целых чисел, пока сумма или произведение не превышают ни модуля $p$ базового поля, ни модуля $2^N$ типа беззнакового целого числа (unsigned integer). При этих ограничениях возможно представить тип \texttt{uN} внутри базового полевого типа, когда $N$ достаточно мало. Однако это небезопасно, и необходимо соблюдать осторожность, чтобы никогда не переполнять ни один из этих модулей и не выходить ниже $0$.
\end{remark}
\begin{example} Чтобы дать интуицию о том, как реальный компилятор (compiler) может преобразовать беззнаковые целые числа (unsigned integers) в переменные (variables) схемы, рассмотрим следующее утверждение на \lgname{PAPER}, которое реализует классический оператор кольцевого сдвига (ring-shift) на типе $u4$ как схему:
\begin{lstlisting}
statement RING_SHIFT{F: F_p, N=4} {
  fn main(x: uN)-> uN {
  	let y : uN ;
  	y <== [x[1],x[2],x[3],x[0]] ;
  	return y ;
  } ;
}
\end{lstlisting}
Во время фазы настройки (setup-phase) компилятор схем (circuit compiler) может затем заменить любое вхождение типа \texttt{uN} на \texttt{N} переменных (variables) булева типа. Используя определение булевых значений, каждая из этих переменных затем преобразуется в полевой тип с булевой системой ограничений. Чтобы увидеть, как это может быть достигнуто, мы используем \lgname{PAPER} как пример и получаем следующую схему:

\begin{center}
\digraph[scale=0.25]{u4SHIFTER}{
  forcelabels=true;
  center=true;
  splines=ortho;
  nodesep= 1.0;

  subgraph clusteru41 {     
    subgraph clusterBOOL1 {
      nb11 [shape=box, label="x_0"] ;
      nb12 [shape=box, label="x_1"] ;    
      nb13 [shape=box, label="x_2"] ;
      nb14 [shape=box, label="x_3"] ;
      color=lightgray;
      label="BOOL[4]" ; 
    }
    color = lightgray ;
    label = "u4"
  }
  
  subgraph clusteru42 {     
    subgraph clusterBOOL2 {
      nb21 [shape=box, label="y_0"] ;
      nb22 [shape=box, label="y_1"] ;    
      nb23 [shape=box, label="y_2"] ;
      nb24 [shape=box, label="y_3"] ;
      color=lightgray;
      label="BOOL[4]" ; 
    }
    color = lightgray ;
    label = "u4"
  }  
  
  subgraph clusterCONS  {
    subgraph clusterBCONS1 {
      nCONSB11 -> nCONSB14 [xlabel="S_1"] ;
      nCONSB11 -> nCONSB16 [xlabel="S_1"] ;
      nCONSB11 -> nCONSB18 [style=dashed, color=lightgrey] ;
      nCONSB12 -> nCONSB15 ;
      nCONSB13 -> nCONSB14 ;
      nCONSB14 -> nCONSB15 ;
      nCONSB15 -> nCONSB16 ;
      nCONSB16 -> nCONSB17 [xlabel="S_2=0  "] ;
      nCONSB11 [shape=box, label="b:BOOL", color=lightgray ] ;
      nCONSB12 [shape=box, label="1"] ;
      nCONSB13 [shape=box, label="-1"] ;
      nCONSB14 [label="*"] ;
      nCONSB15 [label="+"] ;
      nCONSB16 [label="*"] ;
      subgraph clusterBCONS1out {
        nCONSB17 [shape=box, label="0"] ;
        nCONSB18 [shape=box, label="x:F", color=lightgrey] ;
        color = white ; 
      }
      color=lightgray;
      label="x1 : BOOL" ;
    }
    
    subgraph clusterBCONSB2 {
      nCONSB21 -> nCONSB24 [xlabel="S_1"] ;
      nCONSB21 -> nCONSB26 [xlabel="S_1"] ;
      nCONSB21 -> nCONSB28 [style=dashed, color=lightgrey] ;
      nCONSB22 -> nCONSB25 ;
      nCONSB23 -> nCONSB24 ;
      nCONSB24 -> nCONSB25 ;
      nCONSB25 -> nCONSB26 ;
      nCONSB26 -> nCONSB27 [xlabel="S_2=0  "] ;
      nCONSB21 [shape=box, label="b:BOOL", color=lightgray ] ;
      nCONSB22 [shape=box, label="1"] ;
      nCONSB23 [shape=box, label="-1"] ;
      nCONSB24 [label="*"] ;
      nCONSB25 [label="+"] ;
      nCONSB26 [label="*"] ;
      subgraph clusterBCONS1out {
        nCONSB27 [shape=box, label="0"] ;
        nCONSB28 [shape=box, label="x:F", color=lightgrey] ;
        color = white ; 
      }
      color=lightgray;
      label="x2 : BOOL" ;
    }
    
    subgraph clusterBCONSB3 {
      nCONSB31 -> nCONSB34 [xlabel="S_1"] ;
      nCONSB31 -> nCONSB36 [xlabel="S_1"] ;
      nCONSB31 -> nCONSB38 [style=dashed, color=lightgrey] ;
      nCONSB32 -> nCONSB35 ;
      nCONSB33 -> nCONSB34 ;
      nCONSB34 -> nCONSB35 ;
      nCONSB35 -> nCONSB36 ;
      nCONSB36 -> nCONSB37 [xlabel="S_2=0  "] ;
      nCONSB31 [shape=box, label="b:BOOL", color=lightgray ] ;
      nCONSB32 [shape=box, label="1"] ;
      nCONSB33 [shape=box, label="-1"] ;
      nCONSB34 [label="*"] ;
      nCONSB35 [label="+"] ;
      nCONSB36 [label="*"] ;
      subgraph clusterBCONS1out {
        nCONSB37 [shape=box, label="0"] ;
        nCONSB38 [shape=box, label="x:F", color=lightgrey] ;
        color = white ; 
      }
      color=lightgray;
      label="x3 : BOOL" ;
    }
    
    subgraph clusterBCONSB4 {
      nCONSB41 -> nCONSB44 [xlabel="S_1"] ;
      nCONSB41 -> nCONSB46 [xlabel="S_1"] ;
      nCONSB41 -> nCONSB48 [style=dashed, color=lightgrey] ;
      nCONSB42 -> nCONSB45 ;
      nCONSB43 -> nCONSB44 ;
      nCONSB44 -> nCONSB45 ;
      nCONSB45 -> nCONSB46 ;
      nCONSB46 -> nCONSB47 [xlabel="S_2=0  "] ;
      nCONSB41 [shape=box, label="b:BOOL", color=lightgray ] ;
      nCONSB42 [shape=box, label="1"] ;
      nCONSB43 [shape=box, label="-1"] ;
      nCONSB44 [label="*"] ;
      nCONSB45 [label="+"] ;
      nCONSB46 [label="*"] ;
      subgraph clusterBCONS1out {
        nCONSB47 [shape=box, label="0"] ;
        nCONSB48 [shape=box, label="x:F", color=lightgrey] ;
        color = white ; 
      }
      color=lightgray;
      label="x4 : BOOL" ;
    } 
    color = white ; 
  } 
  nb11 -> nCONSB11 ;
  nb12 -> nCONSB21 ;
  nb13 -> nCONSB31 ;
  nb14 -> nCONSB41 ;  
  nCONSB18 -> nb22;
  nCONSB28 -> nb23;
  nCONSB38 -> nb24;
  nCONSB48 -> nb21;
}
\end{center}
Во время фазы доказывающего (prover phase) функция \texttt{main} вызывается с фактическим входом типа \texttt{u4}, скажем, \texttt{x=14}. Доказывающий (Prover) затем должен преобразовать десятичное значение $14$ в его $4$-битное двоичное представление $Bits(14)_2 = <0,1,1,1>$ вне схемы. Затем массив значений поля $x[4] = [0,1,1,1]$ используется как вход в схему. Поскольку все $4$ элемента поля являются либо $0$, либо $1$, четыре булевых ограничения удовлетворены, а выходом является кольцевой сдвиг входного массива из четырёх элементов поля, заданный как $[1,1,1,0]$, что представляет элемент \texttt{u4} $7$.
\end{example}
\paragraph{Операторы беззнаковых целых чисел} Поскольку элементы типа \texttt{uN} представлены как булевы массивы, операторы сдвига реализуются в схемах просто переподключением булевых входных переменных к выходным переменным соответствующим образом.

Логические операторы, такие как \texttt{AND}, \texttt{OR} или \texttt{NOT}, определяются на типе \texttt{uN} путём вызова соответствующих булевых операторов побитово к каждому биту в булевом массиве, представляющем элемент \texttt{uN}.

Сложение и умножение могут быть представлены аналогично тому, как машины представляют эти операции. Сложение может быть реализовано сначала определением схемы \term{full adder} (полного сумматора), а затем объединением $N$ таких схем в схему, которая складывает два элемента из типа \texttt{uN}.
\begin{comment}
To understand the algebraic circuit for the $1$-bit full, recall that we already defined circuits for boolean algebra in the previous section. Abstracting over those circuits, a full adder circuit can then be defined as:
\begin{center}
\digraph[scale=0.4]{ONEBFULLADD}{
  forcelabels=true;
  center=true;
  splines=ortho;
  nodesep= 2.0;
  
  subgraph clusterin {
    nADD01 [shape=box, label="bx_j"] ;
    nADD02 [shape=box, label="by_j"] ;
    nADD03 [shape=box, label="c_(j-1)"] ;
    color = white ;
  }
  
  subgraph clustermid {
    nADD04 [shape=box, label="XOR"] ;
    nADD05 [shape=box, label="XOR"] ;
    nADD06 [shape=box, label="AND"] ;
    nADD07 [shape=box, label="AND"] ;
    nADD08 [shape=box, label="OR"] ;
    
    nADD04 -> {nADD05, nADD06} ;
    nADD06 -> nADD08 ;
    nADD07 -> nADD08 ;
    
    color = white ;
  }
  
  subgraph clusterout {
    nADD09 [shape=box, label="bz_j"] ;
    nADD010 [shape=box, label="c_j"] ;
    color = white ;
  }
  
  nADD01 -> {nADD04, nADD07} ;
  nADD02 -> {nADD04, nADD07} ;
  nADD03 -> {nADD05, nADD06} ;
  nADD05 -> nADD09 ;
  nADD08 -> nADD010 ; 
}
\end{center}
In this circuit the output $bz_j$ is the result of the binary input $bx_j$ and $by_j$, where $bx_j$ is the $j$-th bit of the binary representation of the first summand and $by_j$ is the $j$-th bit of the binary representation of the second summand. The output $c_j$ is the carry bit of the addition and the input $c_{j-1}$ is is the carry bit which is supposed to be either $0$ for $j=0$ or the carry bit output of the previous full adder circuit. Abstracting the $1$-bit adder, we write:
\begin{center}
\digraph[scale=0.6]{BADDMETA}{
  forcelabels=true;
  center=true;
  splines=ortho;
  //nodesep= 2.0;
  n1 [shape=box, label="FULLADD"] ;
  n2 [shape=none, label="bx_j"] ;
  n3 [shape=none, label="by_j"] ;
  n4 [shape=none, label="c_(j-1)"] ;
  n5 [shape=none, label="bz_j"] ;
  n6 [shape=none, label="c_j"] ;
  n2 -> n1 ;
  n3 -> n1 ;
  n4 -> n1 ;
  n1 -> {n5, n6} ;
}
\end{center}
With a circuit definition of the $1$-bit full adder at hand, addition of two uIntN type elements can then be defined as
\begin{center}
\digraph[scale=0.4]{UINTADD}{
  forcelabels=true;
  center=true;
  splines=ortho;
  //nodesep= 2.0;
  
  subgraph clusterin0 {
    nADD01 [shape=box, label="FULLADD"] ;
    nADD02 [shape=none, label="bx_0"] ;
    nADD03 [shape=none, label="by_0"] ;
    nADD05 [shape=none, label="bz_0"] ;
    nADD02 -> nADD01 ;
    nADD03 -> nADD01 ;
    nADD01 -> nADD05 ;
    color = white ;
  }
  
  subgraph clusterin1 {
    nADD11 [shape=box, label="FULLADD"] ;
    nADD12 [shape=none, label="bx_1"] ;
    nADD13 [shape=none, label="by_1"] ;
    //nADD14 [shape=none, label="c_0"] ;
    nADD15 [shape=none, label="bz_1"] ;
    nADD12 -> nADD11 ;
    nADD13 -> nADD11 ;
    //nADD14 -> nADD11 ;
    nADD11 -> nADD15 ;
    color = white ;
  }

  subgraph clusterin2 {
    nADD21 [shape=box, label="FULLADD", color=lightgray] ;
    nADD22 [shape=none, label="bx_j", color=lightgray] ;
    nADD23 [shape=none, label="by_j", color=lightgray] ;
    //nADD24 [shape=none, label="c_(j-1)", color=lightgray] ;
    nADD25 [shape=none, label="bz_j", color=lightgray] ;
    nADD22 -> nADD21 [color=lightgray];
    nADD23 -> nADD21 [color=lightgray];
    //nADD24 -> nADD21 ;
    nADD21 -> nADD25 [color=lightgray] ;
    color = white ;
  }
  
  subgraph clusterinN {
    nADDN1 [shape=box, label="FULLADD"] ;
    nADDN2 [shape=none, label="bx_(N-1)"] ;
    nADDN3 [shape=none, label="by_(N-1)"] ;
    //nADDN4 [shape=none, label="c_(N-2)"] ;
    nADDN5 [shape=none, label="bz_(N-1)"] ;
    nADDN2 -> nADDN1 ;
    nADDN3 -> nADDN1 ;
    //nADDN4 -> nADDN1 ;
    nADDN1 -> nADDN5
    color = white ;
  }
  
  nADD04 [shape=none, label="0"] ;
  nADD04 -> nADD01 ;
  nADD01 -> nADD11 ;
  nADD11 -> nADD21 [style=dashed, color=lightgrey] ;
  nADD21 -> nADDN1  [style=dashed, color=lightgrey] ;
  nADDN6 [shape=none, label="c_out"] ;
  nADDN1 -> nADDN6 ;
  
}
\end{center}
Depending on how the output carry bit is handled we get different definition of addition in this type. One way would be to enforce it to be zero. This way addition in the circuit is only possible if the sum does not exceed $2^N-1$. On the other hand if the carry bit is unconstraint, then the resulting addition is equivalent to modulo $2^N$ arithmetics. Good compilers should therefore always describe explicitly how exactly their implementation of the uintN type behaves such that users don't build their system on false assumptions.

The associated constraint system consists of XXX\sme{add reference} constraints, including the boolean constraints of the representing bits
\paragraph{The boolean Operators} In implementations it is often necessary to execute boolean operations like $ans$, $or$, or $xor$ on elements of the uInt type. Fortunately this easily done by simply applying those operators to every bit separately as shown in XXX\sme{add reference}.  
\end{comment}
% =====================================
\begin{exercise}
Пусть \texttt{F}$=\F_{13}$ и \texttt{N=4} фиксированы, и пусть $x$ имеет тип $uN$. Определите схемы и связанные с ними R1CS для операторов побитового сдвига влево и вправо $x<<2$ и $x>>2$. Выполните связанную схему для $x : u4 = 11$ и сгенерируйте конструктивное доказательство для выполнимости R1CS.
\end{exercise}
\begin{exercise}
Пусть \texttt{F}$=\F_{13}$ и \texttt{N=2} фиксированы. Определите схему и связанную с ней R1CS для оператора сложения $\mathtt{ADD}: uN \times uN \to uN$. Выполните связанную схему, чтобы вычислить $\mathtt{ADD}(2,7)$ для $2,7\in uN$.
\end{exercise}
\begin{exercise} Выполните фазу настройки (setup phase) для следующего кода на \lgname{PAPER} (т.е. ``мозговым компилятором'' скомпилируйте код в схему и выведите связанную с ней R1CS).
\begin{lstlisting}
statement MASK_MERGE {F:F_5, N=4} {
  fn main(pub a : uN, pub b : uN) -> uN {
    let constant mask : uN = 10 ;
    let r : uN ;
    r <== XOR(a,AND(XOR(a,b),mask)) ;
    return r ;
  }
}
\end{lstlisting}
Пусть $L_{mask\_merge}$ будет языком, определённым схемой. Предоставьте конструктивное доказательство знания в $L_{mask\_merge}$ для instance $I=(I_a, I_b) = (14, 7)$.
\end{exercise}
\subsection{Управление потоком выполнения} Большинство языков программирования императивного или функционального стиля имеют некоторое понятие базовых управляющих структур для направления порядка, в котором инструкции вычисляются. Современные компиляторы схем (circuit compilers) обычно предоставляют единственный поток выполнения и предоставляют базовые конструкции потока, реализующие управление потоком в схемах. В этой части мы рассмотрим некоторые базовые конструкции управления потоком и их реализацию в схемах.
\subsubsection{Условное присваивание} При написании высокоуровневого кода, компилируемого в схемы, часто необходимо иметь способ условного присваивания значений или выходных данных вычислений переменным. Один из способов реализовать это во многих языках программирования --- через условный тернарный оператор присваивания $?:$, который ветвит поток управления программы в соответствии с некоторым условием, а затем присваивает выход вычисленной ветви некоторой переменной:
\begin{lstlisting}
variable = condition ? value_if_true : value_if_false  
\end{lstlisting}
В этом описании \texttt{condition} --- булево выражение, а \texttt{value\_if\_true}, а также \texttt{value\_if\_false} --- выражения, вычисляющиеся к тому же типу, что и \texttt{variable}.

В языках программирования, таких как Rust, другой способ записи оператора условного присваивания, более знакомый многим программистам, задаётся как
\begin{lstlisting}
variable = if condition then { 
  value_if_true 
} else { 
  value_if_false 
} 
\end{lstlisting}
В большинстве языков программирования ключевым свойством тернарного оператора присваивания является то, что выражение \texttt{value\_if\_true} вычисляется только если \texttt{condition} вычисляется в true, а выражение \texttt{value\_if\_false} вычисляется только если \texttt{condition} вычисляется в false. На самом деле, компьютерные программы оказались бы очень неэффективными, если бы тернарный оператор вычислял оба выражения независимо от значения \texttt{condition}.

Простой способ реализовать оператор условного присваивания как схему может быть достигнут, если отказаться от требования, что выполняется только одна ветвь условного оператора. Чтобы увидеть это, пусть $b$, $c$ и $d$ --- элементы поля такие, что $b$ булево ограничен. В этом случае следующее уравнение принуждает элемент поля $x$ быть результатом оператора условного присваивания:
\begin{equation}
\label{def:conditional_operator_equation}
x = b\cdot c + (1-b)\cdot d
\end{equation}
Выражая это уравнение через операторы сложения и умножения из~\ref{def:base_field_type}, мы можем ``выровнять'' (flatten)~\ref{def:conditional_operator_equation} в следующую алгебраическую схему:
\begin{center}
\digraph[scale=0.4]{CONDASSIGN}{
	forcelabels=true;
	center=true;
	splines=ortho;
	nodesep= 2.0;

  n1 [shape=box, label="b"]
  n2 [shape=box, label="c"]
  n3 [shape=box, label="d"]
  n4 [shape=box, label="b ? c : d"]
  n5 [shape=box, label="-1"]
  n6 [shape=box, label="1"]
  n7 [label="*"]
  n8 [label="*"]
  n9 [label="*"]
  n10 [label="+"]
  n11 [label="+"]
 
  n1 -> n7 [taillabel= "S_1  "] ;
  n1 -> n8 [taillabel= "S_1 "] ;
  n2 -> n7 [xlabel= "S_2"] ;
  n3 -> n9 [xlabel= "S_4"] ;
  n5 -> n8 ;
  n6 -> n10 ;
  n7 -> n11 [xlabel= "S_3  "] ;
  n8 -> n10 ;
  n9 -> n11 [xlabel= "S_5"] ;
  n10 -> n9 ;
  n11 -> n4 [xlabel= "S_6  "] ;
}
\end{center}
Отметим, что для вычисления корректного присваивания (assignment) к этой схеме необходимы как $S_2$, так и $S_4$. Если входы в узлы $c$ и $d$ сами являются схемами, обе схемы требуют корректных присваиваний (assignments) и поэтому должны быть выполнены. Как следствие, эта реализация оператора условного присваивания должна выполнять все ветви всех схем, что сильно отличается от выполнения обычных компьютерных программ и способствует повышенным вычислительным затратам, которые должен инвестировать любой доказывающий (prover), в отличие от выполнения в других моделях программирования.

Мы можем использовать общую технику из~\ref{sec:circuits_associated_R1CS} для выведения связанной системы ограничений ранга~1 оператора условного присваивания. Мы получаем следующее:
\begin{align*}
S_1 \cdot S_2 & = S_3 \\
(1 - S_1) \cdot S_4 & = S_5 \\
(S_3 + S_5)\cdot 1 &= S_6
\end{align*}
\begin{example} Чтобы дать интуицию о том, как реальный компилятор схем (circuit compiler) может преобразовать любое высокоуровневое описание оператора условного присваивания в схему, рассмотрим следующий код на \lgname{PAPER}:
\begin{lstlisting}
statement CONDITIONAL_OP {F:F_p} {
  fn main(x : F, y : F, b : BOOL) -> F {
    let z : F 
    z <== if b then { 
      ADD(x,y) 
    } else { 
      MUL(x,y) 
    } ;
    return z ; 
  }
}
\end{lstlisting}
``Мозговым компилятором'' (Brain-compiling) скомпилируйте этот код в схему: мы сначала рисуем прямоугольные узлы для всех входных и выходных переменных, а затем преобразуем булев тип в полевой тип вместе со связанным с ним ограничением. Затем мы вычисляем присваивания выходным переменным. Поскольку оператор условного присваивания является функцией верхнего уровня, мы рисуем его схему, а затем рисуем схемы для обоих условных выражений. Мы получаем следующее:
\begin{center}
\digraph[scale=0.4]{CONDITIONALOP}{
  forcelabels=true;
  center=true;
  splines=ortho;
  nodesep= 1.0;

  subgraph clusterinput { 
    nb11 [shape=box, label="x:F"] ;
    nb12 [shape=box, label="y:F"] ;    
    nb13 [shape=box, label="b:BOOL"] ;
    color = lightgray ;
    label = "input"
  }
  
  subgraph clusterout {
    nb21 [shape=box, label="z"] ;
    color=lightgray;
    label="output" ; 
  }  
  
    subgraph clusterBCONS1 {
      nCONSB11 -> nCONSB14 [xlabel="S_1"] ;
      nCONSB11 -> nCONSB16 [xlabel="S_1"] ;
      nCONSB11 -> nCONSB18 [style=dashed, color=lightgrey] ;
      nCONSB12 -> nCONSB15 ;
      nCONSB13 -> nCONSB14 ;
      nCONSB14 -> nCONSB15 ;
      nCONSB15 -> nCONSB16 ;
      nCONSB16 -> nCONSB17 [xlabel="S_2=0  "] ;
      nCONSB11 [shape=box, label="b:BOOL", color=lightgray ] ;
      nCONSB12 [shape=box, label="1"] ;
      nCONSB13 [shape=box, label="-1"] ;
      nCONSB14 [label="*"] ;
      nCONSB15 [label="+"] ;
      nCONSB16 [label="*"] ;
      subgraph clusterBCONS1out {
        nCONSB17 [shape=box, label="0"] ;
        nCONSB18 [shape=box, label="x:F", color=lightgrey] ;
        color = white ; 
      }
      color=lightgray;
      label="x1 : BOOL" ;
    } 
    
    subgraph clusterconditional{
      nCONDI1 [shape=box, label="b", color=lightgrey]
      nCONDI2 [shape=box, label="c", color=lightgrey]
      nCONDI3 [shape=box, label="d", color=lightgrey]
      nCONDI4 [shape=box, label="b ? c : d", color=lightgrey]
      nCONDI5 [shape=box, label="-1"]
      nCONDI6 [shape=box, label="1"]
      nCONDI7 [label="*"]
      nCONDI8 [label="*"]
      nCONDI9 [label="*"]
      nCONDI10 [label="+"]
      nCONDI11 [label="+"]
     
      nCONDI1 -> nCONDI7 [taillabel= "S_1  "] ;
      nCONDI1 -> nCONDI8 [taillabel= "S_1 "] ;
      nCONDI2 -> nCONDI7 [xlabel= "S_2"] ;
      nCONDI3 -> nCONDI9 [xlabel= "S_4"] ;
      nCONDI5 -> nCONDI8 ;
      nCONDI6 -> nCONDI10 ;
      nCONDI7 -> nCONDI11 [xlabel= "S_3  "] ;
      nCONDI8 -> nCONDI10 ;
      nCONDI9 -> nCONDI11 [xlabel= "S_5"] ;
      nCONDI10 -> nCONDI9 ;
      nCONDI11 -> nCONDI4 [xlabel= "S_6  "] ;
      
      color=lightgray;
      label="b ? c : d" ;
    }
    
    subgraph clusterADD {
      nADD1 [shape=box, label="x", color=lightgrey] ;
      nADD2 [shape=box, label="y", color=lightgrey] ;
      nADD3 [label="+"] ;
      nADD4 [shape=box, label="ADD(x,y)", color=lightgrey] ;
      
      nADD1 -> nADD3 ;
      nADD2 -> nADD3 ;
      nADD3 -> nADD4 ;
      
      color=lightgray;
      label="ADD" ;
    }
    
    subgraph clusterMUL {
      nMUL1 [shape=box, label="x", color=lightgrey] ;
      nMUL2 [shape=box, label="y", color=lightgrey] ;
      nMUL3 [label="+"] ;
      nMUL4 [shape=box, label="MUL(x,y)", color=lightgrey] ;
      
      nMUL1 -> nMUL3 ;
      nMUL2 -> nMUL3 ;
      nMUL3 -> nMUL4 ;
      
      color=lightgray;
      label="MUL" ;
    }
    nb11 -> {nMUL1, nADD1} [style=dashed, color=lightgrey] ; 
    nb12 -> {nMUL2, nADD2} [style=dashed, color=lightgrey] ; 
    nb13 -> nCONSB11 [style=dashed, color=lightgrey] ;  
    nCONSB18 -> nCONDI1 [style=dashed, color=lightgrey] ; 
    nADD4 -> nCONDI2 [style=dashed, color=lightgrey] ; 
    nMUL4 -> nCONDI3 [style=dashed, color=lightgrey] ; 
    nCONDI4 -> nb21 [style=dashed, color=lightgrey] ; 
  // outer circuit
}
\end{center}
\end{example}

% NOTE: ZK-Podcast with Alex Özdemir for the proper branching thing in version 2 of the book.

\subsubsection{Циклы} Во многих языках программирования определены различные управляющие структуры циклов, позволяющие разработчикам выполнять выражения с указанным числом повторений. В частности, часто возможно реализовать неограниченные циклы, такие как структура цикла, приведённая ниже:
\begin{lstlisting}
  while true do { }
\end{lstlisting}
Кроме того, часто возможно реализовать структуры циклов, где число шагов выполнения в цикле зависит от входных данных выполнения или промежуточных вычислительных шагов и поэтому неизвестно во время компиляции:
\begin{lstlisting}
  x = 0.5
  while x != 0 do { 
    x = 4*x*(1-x) 
  }
\end{lstlisting}

В противоположность этому, алгебраические схемы и системы ограничений ранга~1 недостаточно общи для выражения произвольных вычислений, а только ограниченных вычислений. Как следствие, невозможно реализовать неограниченные циклы или циклы с границами, неизвестными во время компиляции, в этих моделях. Это легко видно, поскольку схемы ацикличны по определению, а реализация неограниченного цикла как ациклического графа требует схемы неограниченного размера. Однако схемы достаточно общи для выражения ограниченных циклов, где верхняя граница их выполнения известна во время компиляции. Такие циклы могут быть реализованы в схемах путём развёртки (unrolling) цикла.

Как следствие, любой язык программирования, компилируемый в алгебраические схемы, может предоставлять только структуры циклов, где граница является константой, известной во время компиляции. Это означает, что циклы не могут зависеть от входных данных выполнения, а только от параметров времени компиляции.
\begin{example} Чтобы дать интуицию о том, как реальный компилятор схем (circuit compiler) может преобразовать любое высокоуровневое описание ограниченного цикла \texttt{for} в схему, рассмотрим следующий код на \lgname{PAPER}:
\begin{lstlisting}
statement FOR_LOOP {F:F_p, N: unsigned = 4} {
  fn main(fac : F[N]) -> F {
  	let prod[N] : F ;
  	prod[0] <== fac[0] ;
    for unsigned i in 1..N do [{
      prod[i] <== MUL(fac[i], prod[i-1]) ; 
    }
    return prod[N] ;
  }
}
\end{lstlisting}
Отметим, что в такой программе счётчик цикла \texttt{i} не имеет выражения в выведенной схеме. Это высокоуровневый параметр, который сообщает компилятору (compiler), как развернуть (unroll) цикл.

``Мозговым компилятором'' (Brain-compiling) скомпилируйте этот код в схему: мы сначала рисуем прямоугольные узлы для всех входных и выходных переменных, отметив, что счётчик цикла не представлен в схеме. Поскольку все переменные имеют тип \texttt{field}, нам не нужно компилировать никаких ограничений типов. Затем мы вычисляем присваивания выходным переменным, разворачивая (unrolling) цикл в $3$ отдельных оператора присваивания (assignment). Мы получаем:
\begin{center}
\digraph[scale=0.3]{SIMPLELOOP}{
  forcelabels=true;
  center=true;
  splines=ortho;
  //nodesep= 2.0;
  rankdir="LR" ;
  
  subgraph clusterINPUT {
    nIN1 [shape=box, label="fac[0]"] ;
    nIN2 [shape=box, label="fac[1]"] ;
    nIN3 [shape=box, label="fac[2]"] ;
    nIN4 [shape=box, label="fac[3]"] ;
    color = lightgray ;
    label = "input" ;
  }
  
  subgraph clusterOUTPUT {
    nOUT4 [shape=box, label="prod[3]"] ;    
    color = lightgray ;
    label = "output" ;
  }
  
  subgraph clusterMUL1 {
    nMUL11 [shape=box, label="x", color=lightgray] ;
    nMUL12 [shape=box, label="y", color=lightgray] ;
    nMUL13 [label="*"] ;
    nMUL14 [shape=box, label="MUL(x,y)", color=lightgray] ;
    nMUL11 -> nMUL13 ;
    nMUL12 -> nMUL13 ;
    nMUL13 -> nMUL14 ;
    color = lightgray ;
    label="MUL" ;
  }
  
    subgraph clusterMUL2 {
    nMUL21 [shape=box, label="x", color=lightgray] ;
    nMUL22 [shape=box, label="y", color=lightgray] ;
    nMUL23 [label="*"] ;
    nMUL24 [shape=box, label="MUL(x,y)", color=lightgray] ;
    nMUL21 -> nMUL23 ;
    nMUL22 -> nMUL23 ;
    nMUL23 -> nMUL24 ;
    color = lightgray ;
    label="MUL" ;
  }
  
    subgraph clusterMUL3 {
    nMUL31 [shape=box, label="x", color=lightgray] ;
    nMUL32 [shape=box, label="y", color=lightgray] ;
    nMUL33 [label="*"] ;
    nMUL34 [shape=box, label="MUL(x,y)", color=lightgray] ;
    nMUL31 -> nMUL33 ;
    nMUL32 -> nMUL33 ;
    nMUL33 -> nMUL34 ;
    color = lightgray ;
    label="MUL" ;
  }

  nOUT1 [shape=box, label="prod[0]", color=lightgrey] ;
  nOUT2 [shape=box, label="prod[1]", color=lightgrey] ;
  nOUT3 [shape=box, label="prod[2]", color=lightgrey] ;
  
  nIN1 -> nOUT1 [style=dashed, color=lightgrey] ;
  
  nOUT1 -> nMUL12 [style=dashed, color=lightgrey] ;
  nIN2 -> nMUL11 [style=dashed, color=lightgrey] ;
  nMUL14 -> nOUT2 [style=dashed, color=lightgrey] ;
 
  nOUT2 -> nMUL22 [style=dashed, color=lightgrey] ;
  nIN3 -> nMUL21 [style=dashed, color=lightgrey] ;
  nMUL24 -> nOUT3 [style=dashed, color=lightgrey] ; 
  
  nOUT3 -> nMUL32 [style=dashed, color=lightgrey] ;
  nIN4 -> nMUL31 [style=dashed, color=lightgrey] ;
  nMUL34 -> nOUT4 [style=dashed, color=lightgrey] ;  
  
}
\end{center}
\end{example}
\subsection{Двоичные представления полевых элементов} В приложениях часто необходимо обеспечить двоичное представление элементов из полевого типа \texttt{field}. Чтобы вывести подходящую схему над простым полем $\F_p$, пусть $m=|p|_2$ будет наименьшим числом битов, необходимых для представления простого модуля $p$. Тогда битовая строка $<b_0,\ldots,b_{m-1}>\in \{0,1\}^m$ является двоичным представлением элемента поля $x\in\F_p$, тогда и только тогда, когда
\begin{equation}
\label{def:binary_field_rep}
x = b_0\cdot 2^0 + b_1\cdot 2^1 + \ldots + b_{m-1}\cdot 2^{m-1}
\end{equation}
В этом выражении сложение и возведение в степень считаются выполненными в $\F_p$, что корректно определено, поскольку все члены $2^j$ для $0 \leq j \leq m$ являются элементами $\F_p$. Однако, в отличие от двоичного представления беззнаковых целых чисел (unsigned integers) $n\in\N$, это представление не является уникальным в общем случае, поскольку модульный $p$ класс эквивалентности может содержать более одного двоичного представителя.

Учитывая, что лежащее в основе простое поле фиксировано и наиболее значимый бит простого модуля равен $m-1$, следующая схема ``выравнивает'' (flattens) уравнение~\ref{def:binary_field_rep}, предполагая, что все входы $b_0$, $\ldots$, $b_{m-1}$ имеют булев тип.
\begin{center}
\digraph[scale=0.3]{BINARYREP}{
	forcelabels=true;
	center=true;
	splines=ortho;
	nodesep= 2.0;

  subgraph cluster0 {
    n1 [shape=box, label="b_0"] ;
    n2 [shape=box, label="2^0"] ;
    n3 [label="*"] ;

    n1 -> n3 [xlabel="S_0"] ;
    n2 -> n3 ;
    color=white ;
  }

  subgraph cluster1 {
    n4 [shape=box, label="b_1"] ;
    n5 [shape=box, label="2^1"] ;
    n6 [label="*"] ;

    n4 -> n6 [xlabel="S_1  "] ;
    n5 -> n6 ;
    color=white ;
  }

  subgraph cluster2 {
    n7 [shape=box, label="b_2"] ;
    n8 [shape=box, label="2^2"] ;
    n9 [label="*"] ;

    n7 -> n9 [xlabel="S_2  "] ;
    n8 -> n9 ;
    color=white ;
  }

  subgraph cluster3 {
    n10 [shape=box, label="...", color=lightgrey] ;
    color=white ;
  }

  subgraph cluster4 {
    n11 [shape=box, label="b_(m-1)"] ;
    n12 [shape=box, label="2^(m-1)"] ;
    n13 [label="*"] ;

    n11 -> n13 [xlabel="S_(m-1)  "] ;
    n12 -> n13 ;
    color=white ;
  }

  subgraph cluster5 {
    n18 [shape=box, label="x"] ;
    n19 [shape=box, label="-1"] ;
    n20 [label="*"] ;

    n18 -> n20  [xlabel="S_m"] ;
    n19 -> n20 ;
    color=white ;
  }

  n14 [label="+"] ;
  n15 [label="+"] ;
  n16 [label="+", color=lightgrey] ;
  n17 [label="+"] ;
  n21 [label="+"] ;
  n22 [shape=0, label="0"] ;
  n3 -> n14 ;
  n6 -> n14 ; 
  n14 -> n15 ;
  n9 -> n15 ;
  n10 -> n16 [style=dashed, color=lightgrey] ;
  n15 -> n16 [style=dashed, color=lightgrey] ;
  n13 -> n17 ;
  n16 -> n17 [style=dashed, color=lightgrey] ;
  n20 -> n21 ;
  n17 -> n21 ;
  n21 -> n22  [xlabel="W_1=0  "] ;
}
\end{center}
Применяя общее правило преобразования~\ref{sec:circuits_associated_R1CS} для вычисления связанной системы ограничений ранга~1, мы видим, что нам на самом деле нужно только одно ограничение, чтобы обеспечить некоторое двоичное представление любого элемента поля. Мы получаем
$$
(S_0\cdot 2^0 + S_1\cdot 2^1 + \ldots + S_{m-1}\cdot 2^{m-1} -S_m)\cdot 1 = 0
$$
Дан массив \texttt{BOOL[N]} из \texttt{N} булево ограниченных элементов поля и другой элемент поля $x$, схема принуждает \texttt{BOOL[N]} быть одним из двоичных представлений $x$. Если \texttt{BOOL[N]} не является двоичным представлением $x$, не может существовать корректного присваивания (assignment) и, следовательно, решения связанной R1CS.
\begin{example} Рассмотрим простое поле $\F_{13}$. Чтобы вычислить двоичные представления элементов этого поля, мы начинаем с двоичного представления простого модуля $13$, которое равно $Bits(13) = <1,0,1,1>$, поскольку $13= 1\cdot 2^0 + 0\cdot 2^1 + 1\cdot 2^2 + 1\cdot 2^3$. Таким образом, $m=4$ и нам требуется до $4$ битов для представления любого элемента $x\in\F_{13}$.

Чтобы увидеть, что двоичные представления не являются уникальными в общем случае, рассмотрим элемент $2\in \F_{13}$. Он имеет следующие два $4$-битных двоичных представления $Bits(2)=<0,1,0,0>$ и $Bits(2)=<1,1,1,1>$, поскольку в $\F_{13}$ мы имеем
$$
2 = \begin{cases}
0\cdot 2^0 + 1\cdot 2^1 + 0\cdot 2^2 + 0\cdot 2^3\\
1\cdot 2^0 + 1\cdot 2^1 + 1\cdot 2^2 + 1\cdot 2^3
\end{cases}
$$
Это происходит потому, что беззнаковые целые числа (unsigned integers) $2$ и $15$ оба находятся в модульном $13$ классе остатков $2$ и поэтому являются обоими представителями $2$ в $\F_{13}$.

Чтобы увидеть, как работает связанная схема, мы хотим обеспечить двоичное представление $7\in \F_{13}$. Поскольку $m=4$, мы должны обеспечить $4$-битное представление для $7$, которое равно $<1,1,1,0>$, поскольку $7= 1\cdot 2^0 + 1\cdot 2^1 + 1\cdot 2^2 + 0\cdot 2^3$. Корректное присваивание (assignment) схемы поэтому задаётся как $<S_0,S_1,S_2,S_3,S_4>=<1,1,1,0,7>$ и, действительно, присваивание удовлетворяет требуемым $5$ ограничениям, включая $4$ булевых ограничения для $S_0$, $\ldots$, $S_3$:
\begin{align*}
1\cdot (1-1) &= 0 & \text{// булевы ограничения}\\
1\cdot (1-1) &= 0 \\
1\cdot (1-1) &= 0 \\
0\cdot (1-0) &= 0  \\
(1 + 2 + 4 + 0 -7)\cdot 1 &= 0  & \text{// ограничение двоичного представления}
\end{align*}
\end{example}
\subsection{Криптографические примитивы}
В приложениях часто требуется выполнять криптографию в схеме. Для этого базовые криптографические примитивы, такие как хеш-функции или криптография на эллиптических кривых, необходимо реализовывать как схемы. В этом разделе мы приведём несколько базовых примеров того, как реализовать такие примитивы.
\subsubsection{Кривые Твистед-Эдвардса} Реализация криптографии на эллиптических кривых в схемах означает реализацию определяющих уравнений кривых, а также алгебраических операций, таких как групповой закон или скалярное умножение, как схем. Чтобы делать это эффективно, кривая должна быть определена над тем же базовым полем, что и поле, используемое в схеме.

По соображениям эффективности выгодно выбирать эллиптическую кривую так, чтобы все требуемые ограничения и операции могли быть реализованы с как можно меньшим числом вентилей (gates). Кривые Твистед-Эдвардса (Twisted Edwards curves) особенно полезны для этого, поскольку их групповой закон особенно прост, и одно и то же вычисление может использоваться для всех точек кривой, включая точку на бесконечности. Это существенно упрощает схему.
% implementations https://github.com/iden3/circomlib/blob/master/circuits/babyjub.circom
\paragraph{Ограничения кривой Твистед-Эдвардса} Как мы видели в~\ref{sec:edwards}, кривая Твистед-Эдвардса (Twisted Edwards curve) над конечным полем $F$ определяется как множество всех пар точек $(x,y)\in \F\times \F$ таких, что $x$ и $y$ удовлетворяют уравнению $a\cdot x^2+y^2= 1+d\cdot x^2y^2$, и как мы видели в~\examplename{}~\ref{ex:TJJ-circuit_1}, мы можем преобразовать это уравнение в следующую схему:
\begin{center}
\digraph[scale=0.35]{TECCIRCUIT}{
	forcelabels=true;
	center=true;
	splines=ortho;
	nodesep= 2.0;
	n1 -> n4 [label="S_1"labeldistance="4"];
    n1 -> n4 //[ color=lightgray ]
	n2 -> n6 [label="S_2"];
	n2 -> n6 //[ color=lightgray ] ;
	n3 -> n9 /*[ color=lightgray ]*/ ;
	n4 -> n9 [taillabel="S_3", labeldistance="2" /*, color=lightgray */];
	n4 -> n13 [taillabel="S_3", labeldistance="4"];
	n5 -> n10 //[ color=lightgray ] ;
	n6 -> n10 [headlabel="S_4", labeldistance="4"];
	n6 -> n11 [taillabel="S_4", labeldistance="4" /*, color=lightgray */];
	n7 -> n11 /*[ color=lightgray ]*/ ;
	n8 -> n12 /*[ color=lightgray ]*/ ;
	n9 -> n12 /*[ color=lightgray ]*/ ;
	n10 -> n13 //[ color=lightgray ] ; 
	n11 -> n15 /*[ color=lightgray ]*/ ;
	n12 -> n14 /*[ color=lightgray ]*/ ;	
	n13 -> n14 [xlabel="S_5  "  /*, color=lightgray */];
	n14 -> n15 /*[ color=lightgray ]*/ ;
	n15 -> n16 [label="  S_6=0", labeldistance="2" /*, color=lightgray */];
	n1 [shape=box, label="x"];
	n2 [shape=box, label="y"];
	n3 [shape=box, label="-a" /*, color=lightgray */];
	n4 [label="*",  /*, color=lightgray */];
	n5 [shape=box, label="d"];
	n6 [label="*", /*, color=lightgray */];
	n7 [shape=box, label="-1" /*, color=lightgray */];
	n8 [shape=box, label="1" /*, color=lightgray */];
	n9 [label="*" /*, color=lightgray */];
	n10 [label="*"];
	n11 [label="*" /*, color=lightgray */];	
	n12 [label="+" /*, color=lightgray */];	
	n13 [label="*" /*, color=lightgray */];
	n14 [label="+" /*, color=lightgray */];
	n15 [label="+"/*, color=lightgray */];
	n16 [shape=box, label="0" /*, color=lightgray */];		
}
\end{center}
Схема принуждает два входа полевого типа удовлетворять уравнению кривой Твистед-Эдвардса (Twisted Edwards curve), и, как мы знаем из~\examplename{}~\ref{ex:tjj-circuit-2-tjj-R1CS}, связанная система ограничений ранга~1 задаётся как:
\begin{align*}
 S_1 \cdot S_1 &= S_3\\
 S_2 \cdot S_2 &= S_4\\
 (S_4\cdot d)\cdot S_3 &= S_5\\
 (-1\cdot S_4 + S_5 -a\cdot S_3 + 1)\cdot 1 &= 0
\end{align*}
\begin{exercise}
Напишите схему и связанную с ней систему ограничений ранга~1 для кривой Вейерштрасса (Weierstrass curve) данного поля $\F$.
\end{exercise}
\paragraph{Сложение на кривой Твистед-Эдвардса} Как мы видели в~\ref{sec:twisted_ed_group_law}, главное преимущество кривых Твистед-Эдвардса (Twisted Edwards curves) заключается в существовании закона сложения, который не содержит ветвления и верен для всех точек кривой. Более того, нейтральный элемент не задаётся каким-либо вспомогательным символом, а является точкой кривой $(0,1)$. На самом деле, даны две точки $(x_1,y_1)$ и $(x_2,y_2)$ на кривой Твистед-Эдвардса, их сумма определяется как
$$
(x_3,y_3) = \left(\frac{x_1y_2+y_1x_2}{1+d\cdot x_1x_2y_1y_2}, \frac{y_1y_2-a\cdot x_1x_2}{1-d\cdot x_1x_2y_1y_2}\right)
$$
% https://z.cash/technology/jubjub/
Мы можем использовать схему деления из~\ref{def:division_circuit}, чтобы ``выровнять'' (flatten) это уравнение в алгебраическую схему. Входами в схему затем являются не только две точки кривой $(x_1,y_1)$ и $(x_2,y_2)$, но также мультипликативные обратные двух знаменателей $inv_1 = (1+d\cdot x_1x_2y_1y_2)^{-1}$ и $inv_2= (1-d\cdot x_1x_2y_1y_2)^{-1}$, которые любой доказывающий (prover) должен вычислить вне схемы. Мы получаем
\begin{center}
\digraph[scale=0.5]{EDWARDSADD}{
  forcelabels=true;
  center=true;
  splines=ortho;
  //nodesep= 2.0;
  
  subgraph clusterin {
    n1 [shape=box, label="x_1"] ;
    n2 [shape=box, label="x_2"] ;
    n3 [shape=box, label="y_1"] ;
    n4 [shape=box, label="y_2"] ;
      
    n22 [shape=box, label="inv_1"] ;
    n23 [shape=box, label="inv_2"] ;
  
    color=white ;
  }
  
  subgraph clusterout {
    n29 [shape=box, label="x_3"] ;
    n30 [shape=box, label="y_3"] ;
  
    color=white ;
  }

    n5 [shape=box, label="a"] ;
    n6 [shape=box, label="d"] ;
    n7 [shape=box, label="1"] ;
    n8 [shape=box, label="-1"] ;
    
    n9 [label="*"] ; // x_1*y_2
    n10 [label="*"] ; // x_1*x_2
    n11 [label="*"] ; // y_1*x_2
    n12 [label="*"] ; // y_1*y_2
    n13 [label="*"] ; // a*(x_1*x_2)
    n14 [label="*"] ; // -a*(x_1*x_2)
    n15 [label="+"] ; // x_1*y_2 + y_1*x_2
    n16 [label="+"] ; // y_1*y_2 - a*x_1*x_2
    n17 [label="*"] ; // (x_1*x_2)*(y_1*y_2)
    n18 [label="*"] ; // d*(x_1*x_2)*(y_1*y_2)
    n19 [label="*"] ; // -d*(x_1*x_2)*(y_1*y_2)
    n20 [label="+"] ; // 1 + d*(x_1*x_2)*(y_1*y_2)
    n21 [label="+"] ; // 1 - d*(x_1*x_2)*(y_1*y_2)
    
    n24 [label="*"] ; // (1 + d*(x_1*x_2)*(y_1*y_2))*denom_1 =1 
    n25 [label="*"] ; // (1 - d*(x_1*x_2)*(y_1*y_2))*denom_2 =1 
    n26 [shape=box, label="1"] ;
    n27 [label="*"] ; // denom_1*(x_1*y_2 + y_1*x_2) 
    n28 [label="*"] ; // denom_2*(y_1*y_2 - a*x_1*x_2) 
    
    n1 -> n9 [headlabel=" S_1"];
    n1 -> n10 [taillabel="S_1"];
    n2 -> {n10, n11} [taillabel="S_2"];
    n3 -> n11 [headlabel=" S_3"];
    n3 -> n12 [taillabel="S_3"];
    n4 -> n9 [taillabel="S_4"];
    n4 -> n12 [headlabel="  S_4"];
    n5 -> n13 ;
    n6 -> n18 ;
    n7 -> {n20, n21}
    n8 -> {n14, n19} ;
    n9 -> n15 [headlabel=" S_7"] ;
    n10 -> n13 [xlabel="S_8"] ;
    n10 -> n17 [xlabel="S_8"] ;
    n11 -> n15 [xlabel="S_9"] ;
    n12 -> n16 [taillabel="S_10 "] ;  
    n12 -> n17 [xlabel="  S_10"] ;   
    n13 -> n14 ;
    n14 -> n16 ;
    n15 -> n27 ;
    n16 -> n28 ; 
    n17 -> n18 [xlabel="S_11"] ;
    n18 -> {n19, n20} ;
    n19 -> n21 ;
    n20 -> n24 ;
    n21 -> n25 ;
    n22 -> {n24, n27} [xlabel="S_5"] ;
    n23 -> {n25, n28}  [xlabel="S_6"] ;
    n24 -> n26 [xlabel="S_12=1"] ;
    n25 -> n26 [xlabel="S_13=1"] ;
    
    n27 -> n29 [xlabel="S_14"] ;
    n28 -> n30 [xlabel="S_15"] ;
    
}
\end{center}
Используя общую технику из~\ref{sec:circuits_associated_R1CS} для выведения связанной системы ограничений ранга~1, мы получаем следующий результат:
\begin{align*}
S_1 \cdot S_4 & = S_7 \\
S_1 \cdot S_2 & = S_8 \\
S_2 \cdot S_3 & = S_9 \\
S_3 \cdot S_4 & = S_{10} \\
S_8 \cdot S_{10} & = S_{11} \\
S_5 \cdot (1+ d\cdot S_{11}) & = 1 \\
S_6 \cdot (1 - d\cdot S_{11}) & = 1 \\
S_5 \cdot (S_9 + S_7) & = S_{14} \\
S_6 \cdot (S_{10} - a\cdot S_8) & = S_{15}
\end{align*}

\begin{comment}
\begin{example}[Baby-JubJub]
Considering our pen-and-paper baby-jubjub curve over from XXX\sme{add reference}. We recall from XXX\sme{add reference} that $(11,9)$ is a generator for the large prime order subgroup. We therefor already know from XXX\sme{add reference} that
$(11,9) + (7,8) = (11,9) + [3](11,9) = [4](11,9) = (2,9)$. So we execute the circuit  and get
\begin{align*}
x_1 = 11 & = S_1\\
x_2 = 7 & = S_2\\
y_1 = 9 & = S_3 \\
y_2 = 8 & = S_4 \\
1+ 8\cdot S_1\cdot S_2 \cd = 9 & = S_5 \\
y_3 = 8 & = S_6 \\
11 \cdot 8 = 10 & = S_7 \\
11 \cdot 7 =  12 & = S_8 \\
7 \cdot 9 =  11 & = S_9 \\
9 \cdot 8 = 7 & = S_{10} \\
12 \cdot 7 = 6 & = S_{11} \\
S_5 \cdot (1+ d\cdot S_{11}) & = 1 \\
S_6 \cdot (1 - d\cdot S_{11}) & = 1 \\
S_5 \cdot (S_9 + S_7) & = S_{14} \\
S_6 \cdot (S_{10} - a\cdot S_8) & = S_{15}
\end{align*}
\end{example}
\end{comment}
\begin{exercise}
Пусть $\F$ будет полем. Определите схему, которая обеспечивает полевое обращение для точки кривой Твистед-Эдвардса (Twisted Edwards curve) над $\F$.
\end{exercise}
\begin{exercise}
Напишите схему и связанную с ней систему ограничений ранга~1 для закона сложения Вейерштрасса (Weierstrass addition law) данного поля $\F$.
\end{exercise}
\begin{comment}
\paragraph{Twisted Edwards curve inversion} Similar to elliptic curves in Weierstrass form, inversion is cheap on Edwards curves, as the negative of a curve point $-(x,y)$ is given by $(-x,y)$. A curve point $(x_2,y_2)$ is the additive inverse of another curve point $(x_1,y_1)$ precisely if the equation $(x_1,y_1) = (-x_2,y_2)$ holds. We can write this as follows:
$$
\begin{array}{lcl}
x_1 \cdot 1 &=& -x_2 \\
y_1 \cdot 1 &=& y_2
\end{array}
$$
We therefor have a statement of the form $w=(1,x_1,y_1,x_2,y_2)$ and can write the constraints into a matrix equation as
$$
\begin{pmatrix}
0 & 1 & 0 & 0 & 0 \\
0 & 0 & 1 & 0 & 0
\end{pmatrix} \begin{pmatrix} 1 \\ x_1 \\ y_1 \\ x_2 \\ y_2 \end{pmatrix}\odot
\begin{pmatrix}
1 & 0 & 0 & 0 & 0\\
1 & 0 & 0 & 0 & 0
\end{pmatrix} \begin{pmatrix} 1 \\ x_1 \\ y_1 \\ x_2 \\ y_2 \end{pmatrix} =
\begin{pmatrix}
0 & 0 & 0 & -1 & 0\\
0 & 0 & 0 & 0 & 1
\end{pmatrix} \begin{pmatrix} 1 \\ x_1 \\ y_1 \\ x_2 \\ y_2 \end{pmatrix}
$$

In addition we need the following constraints:
$$
\begin{array}{lcl}
x_1 \cdot 1 &=& -x_2 \\
y_1 \cdot 1 &=& y_2
\end{array}
$$
\end{comment}
\begin{comment}
\paragraph{Twisted Edwards curve scalar multiplication} 
% original circuit is here https://iden3-docs.readthedocs.io/en/latest/_downloads/33717d75ab84e11313cc0d8a090b636f/Baby-Jubjub.pdf

Although there are highly optimzed R1CS implementations for scal multiplication on elliptic curves, the basic idea is somewhat simple: Given an elliptic curve $E/\F_r$, a scalar $x\in \F_r$ with binary representation $(b_0,\ldots,b_m)$ and a curve point $P\in E/\F_r$, the scalar multiplication $[x]P$ can be written as
$$
[x]P = [b_0]P + [b_1]([2]P) + [b_2]([4]P) + \ldots + [b_m]([2^m] P)
$$
and since $b_j$ is either $0$ or $1$, $[b_j](kP)$ is either the neutral element of the curve or $[2^j]P$. However, $[2^j]P$ can be computed inductively by curve point doubling, since $[2^j]P= [2]([2^{j-1}]P)$.

So scalar multiplication can be reduced to a loop of length $m$, where the original curve point is repeatedly douled and added to the result, whenever the appropriate bit in the scalar is equal to one.

So to enforce that a curve point $(x_2,y_2)$ is the scalar product $[k](x_1,y_1)$ of a scalar $x\in F_r$ and a curve point $(x_1,y_1)$, we need an R1CS the defines point doubling on the curve (XXX\sme{add reference}) and an R1CS that enforces the binary representation of $x$ (XXX\sme{add reference}). 

In case of a twisted Edwards curve, we can use ordinary addition for doubling, as the constraints works for both cases (doublin is addition, where both arguments are equal). Moreover, $[b](x,y)=(b\cdot x, b\cdot y)$ for boolean $b$. Hence flattening equation XXX\sme{add reference} gives
$$
\begin{array}{lclr}
b_0\cdot x_1 &=& x_{0,1} & // [b_0]P\\
b_0\cdot y_1 &=& y_{0,1}\\

\end{array}
$$
In addition we need to constrain $(b_0,\ldots, b_N)$ to be the binary representation of $x$ and we need to constrain each $b_j$ to be boolean.

As we can see a R1CS for scalar multiplication utilizes many R1CS that we have introduced before. For efficiency and readability it is therefore useful to apply the concept of a gadget (XXX\sme{add reference}). A pseudocode method to derive the associated R1CS could look like this:

%\begin{algorithmic}
%\Require $m$ Bitlength of modulus
%\Statement $w \gets [x,b[m],mid[m]]$
%\State $tmp \gets 0$
%\For{$j\gets 1,\ldots, m$}
%	\State \textbf{Constrain:} $b[j]\cdot (1-b[j]) == 0$
%	\State \textbf{Constrain:} $b[j] \cdot 2^j == mid[j]$
%	\State $tmp = tmp + mid[j]$
%\EndFor
%\State \textbf{Constrain:} $tmp \cdot 1 == x$
%\end{algorithmic}

%\begin{codebox}
%\Procname{$\proc{Insertion-Sort}(A)$}
%\li \For $j \gets 2$ \To $\id{length}[A]$
%\li     \Do$\id{key} \gets A[j]$
%\li         \Comment Insert $A[j]$ into the sorted sequence $A[1 \twodots j-1]$.
%\li         $i \gets j-1$\li         \While $i > 0$ and $A[i] > \id{key}$
%\li             \Do$A[i+1] \gets A[i]$
%\li                 $i \gets i-1$\End
%\li         $A[i+1] \gets \id{key}$\End
%\end{codebox}
\end{comment}
 

